%% ============================================================
%%  احتمال به زبان ساده — جزوهٔ مقدماتی برای دانش‌آموزان دبیرستان
%%  کامپایل: powershell -ExecutionPolicy Bypass -File build.ps1 dabirestan
%%  فصل‌ها در پوشهٔ dabirestan/ هستند.
%% ============================================================
\documentclass[a4paper,11pt,oneside,openany]{book}

\newcommand{\chapterlabel}{فصل}

%% ------------------------------------------------------------
%%  بسته‌ها — همهٔ بسته‌ها باید «قبل از» xepersian بارگذاری شوند
%% ------------------------------------------------------------
\usepackage[a4paper,top=2.6cm,bottom=2.4cm,right=2.3cm,left=2.3cm,
            headheight=16pt,headsep=0.8cm,footskip=1.1cm]{geometry}
\usepackage{amsmath,amssymb,mathtools}
\usepackage{bm}
\usepackage[table]{xcolor}
\usepackage{graphicx}
\usepackage{array,booktabs,tabularx,multirow,multicol}
\usepackage{enumitem}
\usepackage{etoolbox}
\usepackage{calc}
\usepackage{tikz}
\usetikzlibrary{calc,positioning,shapes,arrows.meta,decorations.pathreplacing,
                patterns,shadows,fadings,backgrounds,trees,shapes.geometric}
\usepackage{pgfplots}
\pgfplotsset{compat=1.18}
\usepackage[most]{tcolorbox}
\usepackage{titlesec}
\usepackage{titletoc}
\usepackage{fancyhdr}
\usepackage{setspace}
\usepackage{hyperref}

%% ------------------------------------------------------------
%%  قلم‌ها
%%  قلم فارسی: Estedad (مجوز آزاد OFL) — داخل پوشهٔ fonts/ قرار دارد
%%  تا پروژه روی هر سیستمی (و Overleaf) بدون نصب قلم کامپایل شود.
%%  برای تغییر قلم کافی است فقط همین فایل را ویرایش کنید؛ مثلاً:
%%     \settextfont{XB Niloofar}   یا   \settextfont{Vazirmatn}
%% ------------------------------------------------------------
\usepackage{xepersian}

\settextfont[
  Path=fonts/, Extension=.ttf,
  UprightFont=*-Regular, BoldFont=*-Bold,
  Scale=0.92
]{Estedad}

% قلم‌های کمکی فارسی (برای عنوان‌ها و متن‌های ظریف)
\defpersianfont\headfont[
  Path=fonts/, Extension=.ttf,
  UprightFont=*-ExtraBold, BoldFont=*-Black
]{Estedad}
\defpersianfont\semifont[
  Path=fonts/, Extension=.ttf,
  UprightFont=*-SemiBold, BoldFont=*-Bold, Scale=0.92
]{Estedad}
\defpersianfont\lightfont[
  Path=fonts/, Extension=.ttf,
  UprightFont=*-Light, BoldFont=*-Medium, Scale=0.92
]{Estedad}

% قلم لاتین (متن انگلیسی داخل \lr{...}) — Latin Modern همراه همهٔ توزیع‌ها هست
\setlatintextfont{Latin Modern Roman}

% جداکنندهٔ شماره‌ها (مثلاً «تعریف ۲.۳»)
\SepMark{.}

% فاصلهٔ بین خطوط برای خوانایی بهتر متن فارسی
\linespread{1.32}

%% ------------------------------------------------------------
%%  پالت رنگ جزوه — با تغییر همین‌جا کل ظاهر جزوه عوض می‌شود
%% ------------------------------------------------------------
\definecolor{primary}{HTML}{1B3A6B}     % سرمه‌ای: عنوان‌ها
\definecolor{primarylight}{HTML}{E8EEF7}
\definecolor{accent}{HTML}{E0913A}      % نارنجی گرم: تأکیدها

\definecolor{defcolor}{HTML}{1565C0}    % تعریف: آبی
\definecolor{defbg}{HTML}{EEF4FC}
\definecolor{thmcolor}{HTML}{7A2E8E}    % قضیه/لم/نتیجه: بنفش
\definecolor{thmbg}{HTML}{F6EEF9}
\definecolor{excolor}{HTML}{0F7B6C}     % مثال: سبزآبی
\definecolor{exbg}{HTML}{EAF6F3}
\definecolor{solcolor}{HTML}{4E8B80}
\definecolor{notecolor}{HTML}{C77700}   % نکته: کهربایی
\definecolor{notebg}{HTML}{FFF7E6}
\definecolor{warncolor}{HTML}{C0392B}   % هشدار/اشتباه رایج: قرمز
\definecolor{warnbg}{HTML}{FDEDEB}
\definecolor{histcolor}{HTML}{7B5B45}   % تاریخچه: قهوه‌ای
\definecolor{histbg}{HTML}{F8F2EC}
\definecolor{proofcolor}{HTML}{8A94A6}  % اثبات: خاکستری
\definecolor{sumcolor}{HTML}{1B3A6B}    % خلاصه
\definecolor{sumbg}{HTML}{F1F5FB}
\definecolor{exercolor}{HTML}{455A64}   % تمرین
\definecolor{exerbg}{HTML}{F3F5F6}

% رنگ‌های نمودارها و ون‌ها
\definecolor{setA}{HTML}{3F7FD0}
\definecolor{setB}{HTML}{E0913A}
\definecolor{setC}{HTML}{3AA17E}
\definecolor{vennfill}{HTML}{9BB8E3}
\definecolor{graytext}{HTML}{5F6B7A}

%% ------------------------------------------------------------
%%  صفحه‌آرایی: عنوان جلسه‌ها و بخش‌ها، سرصفحه و پاصفحه، فهرست
%%  نکتهٔ فنی: در محیط راست‌به‌چپ، tikzها را داخل \lr{...} و متن فارسی
%%  داخل گره‌ها را در \rl{...} می‌گذاریم.
%% ------------------------------------------------------------
\setlength{\parindent}{0pt}
\setlength{\parskip}{0.55em}

\hypersetup{
  colorlinks=true, linkcolor=primary, citecolor=defcolor, urlcolor=defcolor,
  pdftitle={Probability and Statistics - Lecture Notes},
  bookmarksnumbered=true
}

\AtBeginDocument{\renewcommand{\contentsname}{فهرست مطالب}}

% برچسب فصل‌ها («جلسهٔ ۸»)
\providecommand{\chapterlabel}{جلسهٔ}

% ---------- عنوان جلسه (chapter) ----------
% xepersian/bidi فرمان‌های chapter را بازتعریف می‌کند؛ پس سرِ فصل را مستقیم می‌سازیم.
\newcommand{\sessionbadge}{%
  \lr{\begin{tikzpicture}[baseline=(b.base)]
    \node[fill=primary, text=white, rounded corners=4pt, inner xsep=13pt, inner ysep=7pt,
          font=\headfont\LARGE](b) {\rl{\chapterlabel{} \thechapter}};
    \fill[accent] ([yshift=-4pt]b.south west) rectangle ([yshift=-7pt]b.south east);
  \end{tikzpicture}}}

\makeatletter
\AtBeginDocument{% بعد از تعریف‌های bidi اجرا شود
\def\@makechapterhead#1{%
  \vspace*{-18pt}%
  {\parindent\z@ \normalfont
   \leavevmode\sessionbadge\par\vspace{12pt}%
   \leavevmode{\color{primary}\headfont\fontsize{23}{36}\selectfont #1\par}%
   \vspace{6pt}\leavevmode\textcolor{primary!20}{\rule{\linewidth}{1.2pt}}\par
   \vspace{14pt}}}
\def\@makeschapterhead#1{%
  \vspace*{-18pt}%
  {\parindent\z@ \normalfont
   \leavevmode{\color{primary}\headfont\fontsize{23}{36}\selectfont #1\par}%
   \vspace{6pt}\leavevmode\textcolor{accent}{\rule{0.22\linewidth}{2.4pt}}\par
   \vspace{18pt}}}

% ---------- فهرست مطالب ----------
\newcommand{\tocbadge}[1]{%
  \lr{\tikz[baseline=(n.base)]\node[fill=primary,text=white,rounded corners=2pt,
      inner xsep=5pt, inner ysep=3pt](n){\rl{\headfont \chapterlabel{} #1}};}\quad}
\renewcommand*\l@chapter[2]{%
  \addpenalty{-\@highpenalty}\addvspace{7pt}%
  {\parindent\z@\leavevmode{\headfont\color{primary}%
   \let\numberline\tocbadge #1\nobreak\hfill #2}\par}\nopagebreak\vspace{3pt}}
}% end AtBeginDocument
\makeatother

% ---------- بخش و زیربخش ----------
\titleformat{\section}
  {\normalfont\headfont\Large\color{primary}}
  {\lr{\tikz[baseline=(n.base)]\node[fill=primary,text=white,rounded corners=2pt,
      inner sep=3.5pt,minimum width=2.1em](n){\rl{\thesection}};}}
  {0.6em}{}
  [\vspace{-3pt}\leavevmode\textcolor{accent}{\rule{\linewidth}{0.9pt}}]
\titlespacing*{\section}{0pt}{20pt}{8pt}

\titleformat{\subsection}
  {\normalfont\headfont\large\color{primary!85!black}}
  {{\color{accent}\thesubsection}}{0.6em}{}
\titlespacing*{\subsection}{0pt}{14pt}{5pt}

\titleformat{\subsubsection}
  {\normalfont\semifont\color{primary!80!black}}
  {}{0em}{{\color{accent}$\blacktriangleleft$}\ }
\titlespacing*{\subsubsection}{0pt}{10pt}{3pt}

\setcounter{secnumdepth}{2}
\setcounter{tocdepth}{1}

% ---------- سرصفحه و پاصفحه ----------
\newcommand{\pagebadge}{\lr{\tikz\node[circle,fill=primary,text=white,inner sep=1.5pt,
                        minimum size=1.8em,font=\small\bfseries]{\rl{\thepage}};}}
\setlength{\headheight}{23pt}      % fancyhdr: حداقل ارتفاع لازم برای سرصفحهٔ فارسی
\addtolength{\topmargin}{-7pt}     % تا جای متن عوض نشود
\pagestyle{fancy}
\fancyhf{}
\renewcommand{\chaptermark}[1]{\markboth{\chapterlabel{} \thechapter:\ #1}{}}
\renewcommand{\sectionmark}[1]{\markright{\thesection\quad #1}}
\fancyhead[R]{\small\color{graytext}\leftmark}
\fancyhead[L]{\small\color{graytext}\coursetitle}
\fancyfoot[C]{\pagebadge}
\renewcommand{\headrulewidth}{0.4pt}
\renewcommand{\headrule}{{\color{primary!30}\hrule height\headrulewidth width\headwidth}}
\fancypagestyle{plain}{%
  \fancyhf{}
  \fancyfoot[C]{\pagebadge}
  \renewcommand{\headrulewidth}{0pt}}

% ---------- فهرست‌ها ----------
\setlist{itemsep=2pt,topsep=4pt,parsep=0pt}
\setlist[itemize,1]{label={\color{accent}\small$\blacktriangleleft$}}
\setlist[itemize,2]{label={\color{primary}\scriptsize$\bullet$}}
\setlist[enumerate,1]{label={\color{primary}\bfseries\arabic*.}}
\setlist[enumerate,2]{label={\color{primary}(\alph*)}}

%% ------------------------------------------------------------
%%  باکس‌ها (tcolorbox)
%%  تعریف، قضیه، لم، گزاره و نتیجه یک شمارندهٔ مشترک دارند
%%  (مثلاً تعریف ۲.۱ ← قضیه ۲.۲ ← نتیجه ۲.۳) تا پیدا کردن‌شان ساده باشد.
%%  مثال‌ها و تمرین‌ها شمارندهٔ جدا دارند.
%%
%%  استفاده:
%%    \begin{definition}{عنوان اختیاری} ... \end{definition}
%%    \begin{theorem}{عنوان}  ... \end{theorem}
%%    \begin{proof} ... \end{proof}
%%    \begin{example}{عنوان} ... \begin{solution} ... \end{solution}\end{example}
%%    \begin{note} ... \end{note}     \begin{warning} ... \end{warning}
%%    \begin{intuition} ... \end{intuition}
%%    \begin{history}{عنوان} ... \end{history}
%%    \begin{keyformula} ... \end{keyformula}
%%  آرگومان اختیاری اول در باکس‌های شماره‌دار برای برچسب است:
%%    \begin{theorem}[label=thm:bayes]{قضیهٔ بیز} ... \end{theorem}
%%  و ارجاع با \ref{thm:bayes}
%%
%%  نکتهٔ فنی (راست‌به‌چپ): هر جا متنِ رنگی یا یک tikz «اولِ پاراگراف» می‌آید،
%%  قبلش \leavevmode لازم است؛ وگرنه رنگ به بقیهٔ پاراگراف نشت می‌کند
%%  یا ترتیب کلمات برعکس می‌شود.
%% ------------------------------------------------------------

% ---------- آیکن کوچک دایره‌ای ----------
\newcommand{\boxicon}[2]{% #1 رنگ, #2 نماد
  \lr{\tikz[baseline=-0.65ex]\node[circle,fill=#1,text=white,inner sep=0pt,
        minimum size=1.3em,font=\footnotesize\bfseries]{#2};}}

% عنوان باکس: «نوع شماره | عنوان»
\newcommand{\boxtitle}[2]{#1 \thetcbcounter\if\relax\detokenize{#2}\relax\else\ \textbar\ #2\fi}

\tcbset{
  basebox/.style={
    enhanced, breakable, arc=3pt, outer arc=3pt,
    left=9pt, right=11pt, top=7pt, bottom=7pt,
    before skip=12pt, after skip=12pt,
    fonttitle=\semifont, pad at break*=3mm,
  },
}

% ---------- تعریف ----------
\newtcolorbox[auto counter, number within=chapter]{definition}[2][]{%
  basebox,
  colback=defbg, colframe=defcolor!70, boxrule=0.6pt,
  borderline east={3.5pt}{0pt}{defcolor},
  top=14pt,
  attach boxed title to top right={xshift=-14pt, yshift*=-\tcboxedtitleheight/2},
  boxed title style={colback=defcolor, colframe=defcolor, arc=2.5pt, boxrule=0pt,
                     left=7pt, right=7pt, top=1.5pt, bottom=1.5pt},
  coltitle=white,
  title={\boxtitle{تعریف}{#2}},
  #1}

% ---------- قضیه و خانواده‌اش ----------
\tcbset{
  thmstyle/.style={
    basebox,
    colback=thmbg, colframe=#1, colbacktitle=#1, coltitle=white,
    boxrule=0.6pt, titlerule=0pt,
    borderline east={3.5pt}{0pt}{#1},
    toptitle=2.5pt, bottomtitle=2.5pt,
  },
}
\newtcolorbox[use counter from=definition]{theorem}[2][]{%
  thmstyle=thmcolor, title={\boxtitle{قضیه}{#2}}, #1}
\newtcolorbox[use counter from=definition]{lemma}[2][]{%
  thmstyle=thmcolor!80!white, title={\boxtitle{لم}{#2}}, #1}
\newtcolorbox[use counter from=definition]{proposition}[2][]{%
  thmstyle=thmcolor!80!white, title={\boxtitle{گزاره}{#2}}, #1}
\newtcolorbox[use counter from=definition]{corollary}[2][]{%
  thmstyle=thmcolor!62!white, title={\boxtitle{نتیجه}{#2}}, #1}

% ---------- اثبات ----------
\newtcolorbox{proof}[1][]{%
  enhanced, breakable, frame hidden, interior hidden,
  borderline east={1.4pt}{0pt}{proofcolor!70},
  left=0pt, right=11pt, top=2pt, bottom=2pt,
  before skip=6pt, after skip=10pt,
  before upper={\leavevmode{\semifont\color{proofcolor!55!black}اثبات.}\ },
  after upper={\hfill{\color{proofcolor}$\blacksquare$}},
  #1}

% ---------- مثال ----------
\newtcolorbox[auto counter, number within=chapter]{example}[2][]{%
  basebox,
  colback=exbg, frame hidden,
  borderline east={3.5pt}{0pt}{excolor},
  sharp corners=east,
  before upper={\leavevmode{\semifont\color{excolor}\boxtitle{مثال}{#2}}\quad},
  #1}

\newenvironment{solution}{%
  \par\vspace{3pt}%
  \noindent\tikz\draw[excolor!45, line width=0.6pt, dash pattern=on 3pt off 2.5pt]
      (0,0) -- (\linewidth,0);\par
  \vspace{1pt}\leavevmode{\semifont\color{solcolor}پاسخ.}\ }{\par}

% ---------- نکته، شهود، هشدار، تاریخچه ----------
\newtcolorbox{note}[1][]{%
  basebox, colback=notebg, frame hidden,
  borderline east={3pt}{0pt}{notecolor}, sharp corners=east,
  before upper={\leavevmode\boxicon{notecolor}{!}\ {\semifont\color{notecolor!85!black}نکته.}\ },
  #1}
\newtcolorbox{intuition}[1][]{%
  basebox, colback=primarylight, frame hidden,
  borderline east={3pt}{0pt}{primary}, sharp corners=east,
  before upper={\leavevmode\boxicon{primary}{?}\ {\semifont\color{primary}شهود.}\ },
  #1}
\newtcolorbox{warning}[1][]{%
  basebox, colback=warnbg, frame hidden,
  borderline east={3pt}{0pt}{warncolor}, sharp corners=east,
  before upper={\leavevmode\boxicon{warncolor}{$\times$}\ {\semifont\color{warncolor}اشتباه رایج.}\ },
  #1}
\newtcolorbox{history}[2][]{%
  basebox, colback=histbg, colframe=histcolor!45, boxrule=0.5pt,
  borderline east={3pt}{0pt}{histcolor},
  before upper={\leavevmode\boxicon{histcolor}{$\mathcal{H}$}\ {\semifont\color{histcolor}گوشهٔ تاریخ%
      \if\relax\detokenize{#2}\relax\else\ \textbar\ #2\fi}\par\smallskip},
  #1}

% ---------- فرمول کلیدی ----------
\newtcolorbox{keyformula}[1][]{%
  enhanced, colback=white, colframe=accent, boxrule=1pt, arc=4pt,
  left=12pt, right=12pt, top=0pt, bottom=4pt,
  before skip=10pt, after skip=10pt,
  width=0.86\linewidth, center,
  #1}

% ---------- کارت شروع جلسه ----------
% \sessioncard{تاریخ}{موضوع طبق سیلابس}{منابع}{اهداف (به صورت itemize)}
\newcommand{\sessioncard}[4]{%
  \begin{tcolorbox}[enhanced, colback=primarylight, frame hidden, arc=4pt,
      left=12pt, right=12pt, top=9pt, bottom=9pt, before skip=4pt, after skip=16pt,
      borderline north={2.5pt}{0pt}{accent}]
    \begin{tabularx}{\linewidth}{@{}>{\semifont\color{primary}}l X@{}}
      تاریخ جلسه: & #1\\[3pt]
      موضوع در سیلابس: & #2\\[3pt]
      منابع این جلسه: & #3
    \end{tabularx}
    \par\medskip
    \leavevmode{\semifont\color{primary}در پایان این جلسه باید بتوانید:}
    #4
  \end{tcolorbox}}

% ---------- خلاصهٔ جلسه ----------
\newtcolorbox{summary}[1][]{%
  enhanced, breakable, colback=sumbg, colframe=sumcolor, boxrule=0.8pt, arc=4pt,
  left=11pt, right=11pt, top=8pt, bottom=8pt, before skip=14pt,
  colbacktitle=sumcolor, coltitle=white, fonttitle=\headfont\large,
  title={\leavevmode\boxicon{accent}{$\Sigma$}\ \ خلاصهٔ جلسه}, #1}

% ---------- تمرین ----------
\newcounter{exercise}[chapter]
\renewcommand{\theexercise}{\thechapter.\arabic{exercise}}
\newenvironment{exercise}[1][]{%
  \refstepcounter{exercise}\par\medskip
  \noindent\leavevmode{\semifont\color{exercolor}\tikz[baseline=(n.base)]\node[fill=exercolor!10,
     draw=exercolor!45, rounded corners=2pt, inner sep=2.5pt](n)
     {\rl{تمرین \theexercise}};}%
  \if\relax\detokenize{#1}\relax\else\ {\color{graytext}(#1)}\fi\ }{\par}
\newcommand{\hard}{{\color{warncolor}$\bigstar$}}   % نشانهٔ تمرین دشوارتر
\newcommand{\answer}[1]{\par\leavevmode{\small\color{graytext}{\semifont جواب:} #1}}

% ---------- زیرنویس شکل (شماره‌گذاری خودکار با شمارندهٔ figure) ----------
% استفاده: \figcaption[fig:label]{متن زیرنویس}   و ارجاع با \ref{fig:label}
\newcommand{\figcaption}[2][]{%
  \par\refstepcounter{figure}%
  \if\relax\detokenize{#1}\relax\else\label{#1}\fi
  \leavevmode{\small\color{graytext}{\semifont شکل \thefigure:}\ #2}\par}

% ---------- بیشتر بدانید (مطالب تکمیلی، خارج از امتحان) ----------
% استفاده: \begin{moreinfo}{عنوان} ... \end{moreinfo}
\definecolor{morecolor}{HTML}{2E7D8C}
\definecolor{morebg}{HTML}{EDF6F8}
\newtcolorbox{moreinfo}[2][]{%
  basebox, colback=morebg, colframe=morecolor!55, boxrule=0.6pt,
  borderline east={3.5pt}{0pt}{morecolor},
  top=14pt,
  attach boxed title to top right={xshift=-14pt, yshift*=-\tcboxedtitleheight/2},
  boxed title style={colback=morecolor, colframe=morecolor, arc=2.5pt, boxrule=0pt,
                     left=7pt, right=7pt, top=1.5pt, bottom=1.5pt},
  coltitle=white,
  title={بیشتر بدانید\if\relax\detokenize{#2}\relax\else\ \textbar\ #2\fi},
  #1}

% محیط شکل: شکل و زیرنویسش را در یک صفحه نگه می‌دارد (بدون شناور شدن)
\newenvironment{figblock}{\par\medskip\noindent\begin{minipage}{\linewidth}\centering}{\end{minipage}\par\medskip}

%% ------------------------------------------------------------
%%  دستورهای کمکی
%% ------------------------------------------------------------
% اصطلاح جدید:  \term{فضای نمونه}{Sample Space}
\newcommand{\term}[2]{{\semifont\color{primary}#1}\,{\small(\lr{#2})}}
% تأکید در متن فارسی (قلم فارسی ایتالیک ندارد؛ به جای آن رنگ و وزن)
\renewcommand{\emph}[1]{{\semifont\color{primary!85!black}#1}}
\newcommand{\key}[1]{{\semifont\color{accent!80!black}#1}}
\newcommand{\en}[1]{\lr{#1}}

% ---------- ریاضی ----------
\renewcommand{\P}{\mathbb{P}}
\newcommand{\R}{\mathbb{R}}
\newcommand{\N}{\mathbb{N}}
\newcommand{\Z}{\mathbb{Z}}
\newcommand{\Q}{\mathbb{Q}}
\newcommand{\F}{\mathcal{F}}
\newcommand{\set}[1]{\left\{#1\right\}}
\newcommand{\setc}[2]{\left\{#1 \,:\, #2\right\}}
\newcommand{\card}[1]{\left|#1\right|}
\newcommand{\given}{\,\middle|\,}
\newcommand{\cond}[2]{\P\left(#1 \given #2\right)}
\newcommand{\ind}{\mathbf{1}}
\newcommand{\defeq}{\coloneqq}
% عبارت فارسی داخل فرمول، با ترتیب درست کلمات:  $\P(\fa{دقیقاً $k$ شیر})$
% (فارسی و ریاضی را به شکل x\text{...}y در هم نیامیزید؛ ترتیب کلمات برعکس می‌شود)
\newcommand{\fa}[1]{\text{\rl{#1}}}


\renewcommand{\P}{P}                 % در دبیرستان احتمال را با P معمولی می‌نویسیم
\newcommand{\nS}{n(S)}

\newcommand{\coursetitle}{احتمال به زبان ساده}
\hypersetup{pdftitle={Probability for High School}}

% کادر «در این فصل یاد می‌گیرید»
\newcommand{\chapterintro}[1]{%
  \begin{tcolorbox}[enhanced, colback=primarylight, frame hidden, arc=4pt,
      left=12pt, right=12pt, top=9pt, bottom=9pt, before skip=4pt, after skip=16pt,
      borderline north={2.5pt}{0pt}{accent}]
    \leavevmode{\semifont\color{primary}در این فصل یاد می‌گیرید:}
    #1
  \end{tcolorbox}}

\begin{document}

%% ------------------------- جلد -------------------------
\begin{titlepage}
\thispagestyle{empty}
\begin{tikzpicture}[remember picture, overlay]
  \fill[primary] (current page.north west) rectangle ([yshift=-13cm]current page.north east);
  \fill[accent] ([yshift=-13cm]current page.north west) rectangle ([yshift=-13.3cm]current page.north east);
  % تاس‌ها و سکه‌ها
  \begin{scope}[shift={([xshift=3.2cm,yshift=-3.4cm]current page.north west)}, rotate=14]
    \draw[white, line width=1.4pt, rounded corners=6pt, fill=white, fill opacity=0.08] (0,0) rectangle (2,2);
    \foreach \p in {(0.5,0.5),(1.5,1.5),(1,1),(0.5,1.5),(1.5,0.5)} \fill[white] \p circle (0.15);
  \end{scope}
  \begin{scope}[shift={([xshift=6.2cm,yshift=-5.4cm]current page.north west)}, rotate=-12]
    \draw[white, line width=1.4pt, rounded corners=6pt, fill=white, fill opacity=0.08] (0,0) rectangle (1.6,1.6);
    \foreach \p in {(0.4,0.4),(1.2,1.2)} \fill[white] \p circle (0.13);
  \end{scope}
  \foreach \x/\y/\r in {3.5/-8.6/0.75, 5.6/-9.9/0.55, 2.4/-10.6/0.45} {
    \draw[accent!60!white, line width=1.5pt, fill=accent, fill opacity=0.25] ([xshift=\x cm,yshift=\y cm]current page.north west) circle (\r);
  }
  % عنوان
  \node[anchor=north east, text=white] at ([xshift=-2cm, yshift=-4.2cm]current page.north east)
    {\rl{\headfont\fontsize{40}{54}\selectfont احتمال}};
  \node[anchor=north east, text=accent!40!white] at ([xshift=-2cm, yshift=-6.3cm]current page.north east)
    {\rl{\headfont\fontsize{30}{40}\selectfont به زبان ساده}};
  \node[anchor=north east, text=white!80!primary] at ([xshift=-2cm, yshift=-8.6cm]current page.north east)
    {\rl{\fontsize{15}{22}\selectfont جزوهٔ آموزشی برای دانش‌آموزان دبیرستان}};
  % فهرست کوتاه
  \node[anchor=north east, text=primary] at ([xshift=-2cm, yshift=-15cm]current page.north east)
    {\rl{\headfont\Large در این جزوه}};
  \foreach \t [count=\i] in {{احتمال چیست؟ آزمایش، فضای نمونه و پیشامد},
                             {شمارش: اصل ضرب، جایگشت و ترکیب},
                             {احتمال در فضاهای هم‌شانس و قوانین احتمال},
                             {احتمال شرطی و نمودار درختی},
                             {پیشامدهای مستقل},
                             {قانون احتمال کل و قضیهٔ بیز}} {
    \node[anchor=east, fill=primarylight, text=primary, rounded corners=3pt, minimum height=0.85cm,
          inner xsep=10pt, font=\small] at ([xshift=-2cm, yshift={-15.9cm-1.02cm*\i}]current page.north east)
      {\rl{{\headfont\color{accent}\chapterlabel{} \i}\quad \t}};
  }
\end{tikzpicture}
\end{titlepage}

\frontmatter
\tableofcontents

\chapter*{به دانش‌آموز}
\addcontentsline{toc}{chapter}{به دانش‌آموز}
\markboth{به دانش‌آموز}{}

هر روز با جمله‌هایی مثل «احتمالاً فردا باران می‌آید»، «شانس برد این تیم کم است» یا «این بازی شانسی است» روبه‌رو می‌شویم. علم احتمال به ما یاد می‌دهد این «شانس»ها را \emph{با عدد} بیان کنیم و درست درباره‌شان فکر کنیم. در این جزوه از ساده‌ترین مفهوم‌ها شروع می‌کنیم و قدم‌به‌قدم تا قضیهٔ بیز پیش می‌رویم؛ قضیه‌ای که امروز در پزشکی، هوش مصنوعی و فیلترهای ایمیل استفاده می‌شود.

\textbf{چطور از این جزوه استفاده کنید؟}
\begin{itemize}
  \item هر فصل با کادر «در این فصل یاد می‌گیرید» شروع می‌شود و با «خلاصه» و «تمرین» تمام می‌شود.
  \item \textbf{مثال‌ها را خودتان حل کنید:} اول صورت مثال را بخوانید، کمی فکر کنید و بعد سراغ پاسخ بروید.
  \item کادرهای رنگی کمکتان می‌کنند:
\end{itemize}

\begin{definition}[title={تعریف}]{}
تعریف‌های اصلی در کادر آبی هستند.
\end{definition}
\begin{theorem}[title={قانون}]{}
قانون‌ها و فرمول‌های مهم در کادر بنفش هستند.
\end{theorem}
\begin{note}نکته‌های مهم.\end{note}
\begin{warning}اشتباه‌هایی که خیلی‌ها مرتکب می‌شوند.\end{warning}
\begin{moreinfo}{}مطالب اضافه برای کنجکاوها؛ برای امتحان لازم نیست.\end{moreinfo}

برای بیشتر تمرین‌ها جواب نهایی آمده است تا بتوانید کارتان را بررسی کنید. موفق باشید!

\mainmatter
% !TEX root = ../dabirestan.tex
%% ============================================================
%%  فصل ۱ — احتمال چیست؟ آزمایش، فضای نمونه و پیشامد
%% ============================================================
\chapter{احتمال چیست؟}
\label{hs:ch1}

\chapterintro{\begin{itemize}
  \item احتمال یعنی چه و چرا آن را با عددی بین ۰ و ۱ نشان می‌دهیم؛
  \item آزمایش تصادفی، فضای نمونه و پیشامد چیست؛
  \item چطور پیشامدها را با «یا»، «و» و «نه» ترکیب کنیم؛
  \item احتمال تجربی را چطور با تکرار آزمایش به دست آوریم.
\end{itemize}}

%% ------------------------------------------------------------
\section{شانس را با عدد بیان کنیم}
%% ------------------------------------------------------------

در زندگی روزمره مدام دربارهٔ «شانس» حرف می‌زنیم: «احتمالاً امروز باران نمی‌آید»، «شانس قبولی‌ام زیاد است»، «برنده شدن در قرعه‌کشی تقریباً غیرممکن است». در ریاضی این شانس‌ها را با یک \emph{عدد بین ۰ و ۱} بیان می‌کنیم و اسمش را \term{احتمال}{probability} می‌گذاریم:

\begin{figblock}
\begin{tikzpicture}[x=13cm]
  \shade[left color=warncolor!35, right color=excolor!40] (0,-0.18) rectangle (1,0.18);
  \draw[primary, thick] (0,-0.18) rectangle (1,0.18);
  \foreach \x/\l/\t in {0/0/{غیرممکن}, 0.25/{\tfrac14}/{بعید}, 0.5/{\tfrac12}/{پنجاه‌پنجاه}, 0.75/{\tfrac34}/{محتمل}, 1/1/{حتمی}} {
    \draw[primary, thick] (\x,-0.28) -- (\x,0.28);
    \node[font=\small] at (\x,-0.62) {\lr{$\l$}};
    \node[font=\small\semifont, text=primary] at (\x,0.65) {\rl{\t}};
  }
  \node[font=\scriptsize, text=graytext, align=center] at (0,1.25) {\rl{خورشید از غرب طلوع کند}};
  \node[font=\scriptsize, text=graytext, align=center] at (0.5,1.25) {\rl{سکه شیر بیاید}};
  \node[font=\scriptsize, text=graytext, align=center] at (1,1.25) {\rl{فردا صبح شود}};
\end{tikzpicture}

\figcaption[hs:fig:scale]{مقیاس احتمال: هر چه احتمال به ۱ نزدیک‌تر باشد، رخ دادن آن پیشامد محتمل‌تر است.}
\end{figblock}

\begin{itemize}
  \item احتمال \textbf{۰} یعنی آن اتفاق \emph{هرگز} نمی‌افتد.
  \item احتمال \textbf{۱} یعنی آن اتفاق \emph{حتماً} می‌افتد.
  \item احتمال $\tfrac12$ یعنی شانس افتادن و نیفتادنش برابر است، مثل شیر آمدن یک سکهٔ سالم.
\end{itemize}

احتمال را گاهی به صورت درصد هم می‌گوییم: احتمال $0.3$ همان ۳۰٪ است. پیش‌بینی هواشناسی «۷۰٪ احتمال بارش» یعنی احتمال باران $0.7$ است.

%% ------------------------------------------------------------
\section{آزمایش تصادفی و فضای نمونه}
%% ------------------------------------------------------------

\begin{definition}[label=hs:def:experiment]{آزمایش تصادفی، برآمد، فضای نمونه}
\begin{itemize}
  \item \term{آزمایش تصادفی}{random experiment} کاری است که نتیجه‌اش را از قبل نمی‌دانیم؛ مثل پرتاب سکه یا تاس.
  \item هر نتیجهٔ ممکن آزمایش را یک \term{برآمد}{outcome} می‌نامیم.
  \item مجموعهٔ \emph{همهٔ} برآمدهای ممکن را \term{فضای نمونه}{sample space} می‌نامیم و با $S$ نشان می‌دهیم. تعداد اعضای $S$ را $n(S)$ می‌نویسیم.
\end{itemize}
\end{definition}

\begin{example}[label=hs:ex:spaces]{چند فضای نمونه}
\begin{enumerate}
  \item پرتاب یک سکه: دو برآمد «شیر» و «خط» داریم. برای کوتاه‌نویسی شیر را $H$ و خط را $T$ می‌نامیم: \ $S=\{H,T\}$ و $n(S)=2$.
  \item پرتاب یک تاس: \ $S=\{1,2,3,4,5,6\}$ و $n(S)=6$.
  \item پرتاب دو سکه (یکی پس از دیگری): \ $S=\{HH,\ HT,\ TH,\ TT\}$ و $n(S)=4$.
  \item انتخاب یک روز هفته به تصادف: \ $n(S)=7$.
\end{enumerate}
\end{example}

برای آزمایش‌هایی که چند مرحله دارند، \emph{نمودار درختی} کمک می‌کند هیچ برآمدی را از قلم نیندازیم:

\begin{figblock}
\begin{tikzpicture}[
  grow=left, level distance=2.6cm,
  level 1/.style={sibling distance=2.2cm},
  level 2/.style={sibling distance=1cm},
  every node/.style={font=\small},
  st/.style={draw=primary!60, fill=primarylight, rounded corners=2pt, inner sep=3pt},
  edge from parent/.style={draw=primary!55, thick}]
  \node[circle, fill=primary, inner sep=2pt] {}
    child { node[st] {\lr{$H$}}
      child { node[st] (hh) {\lr{$H$}} }
      child { node[st] (ht) {\lr{$T$}} } }
    child { node[st] {\lr{$T$}}
      child { node[st] (th) {\lr{$H$}} }
      child { node[st] (tt) {\lr{$T$}} } };
  \foreach \n/\l in {hh/HH, ht/HT, th/TH, tt/TT} \node[anchor=east, font=\small, text=accent!80!black] at ([xshift=-3mm]\n.west) {\lr{$\l$}};
  \node[font=\scriptsize, text=graytext] at (-2.6,2.0) {\rl{سکهٔ اول}};
  \node[font=\scriptsize, text=graytext] at (-5.2,2.0) {\rl{سکهٔ دوم}};
\end{tikzpicture}

\figcaption[hs:fig:tree2]{نمودار درختی پرتاب دو سکه: هر مسیر از ریشه تا انتها یک برآمد است.}
\end{figblock}

\begin{warning}
در پرتاب دو سکه، برآمدهای $HT$ (اولی شیر، دومی خط) و $TH$ (اولی خط، دومی شیر) \emph{دو برآمد متفاوت}‌اند. اگر فکر کنید فقط سه حالت «دو شیر، دو خط، یکی از هر کدام» وجود دارد، در محاسبهٔ احتمال اشتباه می‌کنید (در فصل ۳ دلیلش را می‌بینیم).
\end{warning}

%% ------------------------------------------------------------
\section{پیشامد}
%% ------------------------------------------------------------

\begin{definition}[label=hs:def:event]{پیشامد}
هر زیرمجموعه از فضای نمونه یک \term{پیشامد}{event} نام دارد. اگر برآمد آزمایش عضوی از پیشامد $A$ باشد، می‌گوییم \emph{$A$ رخ داده است}.
\end{definition}

\begin{example}[label=hs:ex:events]{پیشامدهایی در پرتاب تاس}
در پرتاب یک تاس ($S=\{1,2,3,4,5,6\}$):
\begin{itemize}
  \item «عدد زوج بیاید»: \ $A=\{2,4,6\}$.
  \item «عدد بیشتر از ۴ بیاید»: \ $B=\{5,6\}$.
  \item «عدد ۷ بیاید»: \ $\varnothing$ (مجموعهٔ تهی)؛ این پیشامد \emph{غیرممکن} است.
  \item «عددی بین ۱ تا ۶ بیاید»: \ $S$؛ این پیشامد \emph{حتمی} است.
\end{itemize}
اگر تاس ۶ بیاید، هم $A$ و هم $B$ رخ داده‌اند. اگر ۲ بیاید، $A$ رخ داده ولی $B$ نه.
\end{example}

\subsection{ترکیب پیشامدها: «یا»، «و»، «نه»}

پیشامدها مجموعه‌اند، پس با عمل‌های مجموعه‌ای ترکیبشان می‌کنیم:

\begin{center}
\renewcommand{\arraystretch}{1.4}
\begin{tabular}{>{\semifont}l c l}
\toprule
به زبان عادی & نماد & یعنی \\
\midrule
«$A$ \emph{یا} $B$» (دست‌کم یکی) & $A\cup B$ & اجتماع \\
«$A$ \emph{و} $B$» (هر دو) & $A\cap B$ & اشتراک \\
«$A$ رخ \emph{ندهد}» & $A'$ & متمم (همهٔ برآمدهایی که در $A$ نیستند)\\
\bottomrule
\end{tabular}
\end{center}

\begin{figblock}
\begin{tikzpicture}[scale=0.62]
  \def\univ{(-1.7,-1.35) rectangle (2.9,1.35)}
  \def\sA{(0,0) circle (1)}
  \def\sB{(1.2,0) circle (1)}
  \tikzset{venn/.style={draw=primary, thick}}
  \begin{scope}[xshift=10.8cm]
    \fill[vennfill] \sA; \fill[vennfill] \sB; \draw[venn] \sA \sB; \draw[graytext] \univ;
    \node at (-0.4,0.25) {\lr{$A$}}; \node at (1.6,0.25) {\lr{$B$}};
    \node[font=\small] at (0.6,-1.85) {\lr{$A\cup B$}};
  \end{scope}
  \begin{scope}[xshift=5.4cm]
    \begin{scope}\clip \sA; \fill[vennfill] \sB; \end{scope}
    \draw[venn] \sA \sB; \draw[graytext] \univ;
    \node at (-0.4,0.25) {\lr{$A$}}; \node at (1.6,0.25) {\lr{$B$}};
    \node[font=\small] at (0.6,-1.85) {\lr{$A\cap B$}};
  \end{scope}
  \begin{scope}
    \fill[vennfill] \univ; \fill[white] \sA; \draw[venn] \sA; \draw[graytext] \univ;
    \node at (0,0) {\lr{$A$}};
    \node[font=\small] at (0.6,-1.85) {\lr{$A'$}};
  \end{scope}
\end{tikzpicture}

\figcaption[hs:fig:venn]{نمودار وِن: مستطیل فضای نمونه $S$ است و ناحیهٔ رنگی، پیشامد حاصل.}
\end{figblock}

\begin{definition}[label=hs:def:disjoint]{پیشامدهای ناسازگار}
دو پیشامد $A$ و $B$ را \term{ناسازگار}{mutually exclusive} می‌گوییم اگر نتوانند با هم رخ دهند؛ یعنی $A\cap B=\varnothing$. مثلاً در پرتاب تاس، «زوج آمدن» و «فرد آمدن» ناسازگارند.
\end{definition}

\begin{example}[label=hs:ex:combine]{ترکیب پیشامدها}
در پرتاب تاس، $A=\{2,4,6\}$ (زوج) و $B=\{4,5,6\}$ (بیشتر از ۳). پیشامدهای $A\cup B$، $A\cap B$ و $A'$ را بنویسید و به زبان عادی بیان کنید.
\begin{solution}
\begin{itemize}
  \item $A\cup B=\{2,4,5,6\}$: «زوج یا بیشتر از ۳».
  \item $A\cap B=\{4,6\}$: «زوج و بیشتر از ۳».
  \item $A'=\{1,3,5\}$: «زوج نباشد»، یعنی «فرد».
\end{itemize}
\end{solution}
\end{example}

%% ------------------------------------------------------------
\section{احتمال تجربی: با تکرار آزمایش}
%% ------------------------------------------------------------

یک راه پیدا کردن احتمال این است که آزمایش را \emph{بارها} تکرار کنیم و ببینیم پیشامد مورد نظر چند بار رخ می‌دهد.

\begin{definition}[label=hs:def:empirical]{احتمال تجربی (فراوانی نسبی)}
اگر آزمایشی را $N$ بار تکرار کنیم و پیشامد $A$ در $N_A$ بار رخ دهد، کسر $\dfrac{N_A}{N}$ را \term{فراوانی نسبی}{relative frequency} یا \emph{احتمال تجربی} $A$ می‌نامیم.
\end{definition}

\begin{example}[label=hs:ex:coins-history]{آزمایش‌های واقعی پرتاب سکه}
چند دانشمند در گذشته واقعاً سکه را هزاران بار پرتاب کرده‌اند:
\begin{center}
\renewcommand{\arraystretch}{1.3}
\begin{tabular}{lccc}
\toprule
آزمایشگر & تعداد پرتاب & تعداد شیر & فراوانی نسبی \\
\midrule
بوفون & $4040$ & $2048$ & $0.507$ \\
پیرسون & $24000$ & $12012$ & $0.5005$ \\
\bottomrule
\end{tabular}
\end{center}
هر چه تعداد پرتاب‌ها بیشتر شود، فراوانی نسبی به $\tfrac12$ نزدیک‌تر می‌شود. این واقعیت را \emph{قانون اعداد بزرگ} می‌نامند.
\end{example}

\begin{example}[label=hs:ex:pin]{پونز}
وقتی یک پونز را روی میز می‌اندازیم، یا «نوک به بالا» می‌افتد یا «به پهلو». احتمال «نوک به بالا» چقدر است؟
\begin{solution}
اینجا دو حالت \emph{هم‌شانس نیستند} (شکل پونز متقارن نیست)، پس نمی‌توانیم بگوییم $\tfrac12$. تنها راه، آزمایش است: مثلاً اگر پونز را ۲۰۰ بار بیندازیم و ۱۲۴ بار نوک به بالا بیفتد، احتمال تجربی $\tfrac{124}{200}=0.62$ است. با تکرار بیشتر، تخمین دقیق‌تری به دست می‌آید.
\end{solution}
\end{example}

\begin{warning}
\textbf{خطای قمارباز:} اگر سکهٔ سالمی ۵ بار پشت سر هم شیر بیاید، احتمال اینکه دفعهٔ ششم خط بیاید همچنان $\tfrac12$ است، نه بیشتر! سکه «حافظه» ندارد و نمی‌داند قبلاً چه آمده است. قانون اعداد بزرگ دربارهٔ تعداد \emph{خیلی زیاد} پرتاب است، نه اینکه نتیجه‌ها همدیگر را «جبران» کنند.
\end{warning}

در فصل ۳ راه دوم را می‌بینیم: وقتی برآمدها \emph{هم‌شانس} باشند، احتمال را بدون آزمایش و فقط با \emph{شمردن} حساب می‌کنیم. برای همین، فصل بعد را به شمارش اختصاص داده‌ایم.

%% ------------------------------------------------------------
\section*{خلاصهٔ فصل ۱}
\addcontentsline{toc}{section}{خلاصه و تمرین‌ها}
%% ------------------------------------------------------------

\begin{summary}[title={\leavevmode\boxicon{accent}{$\Sigma$}\ \ خلاصهٔ فصل}]
\begin{itemize}
  \item احتمال عددی بین ۰ (غیرممکن) و ۱ (حتمی) است.
  \item \textbf{فضای نمونه} $S$ مجموعهٔ همهٔ برآمدهای ممکن است و هر زیرمجموعهٔ آن یک \textbf{پیشامد} است.
  \item «یا» $=\cup$، \ «و» $=\cap$، \ «نه» $=$ متمم $A'$. پیشامدهای \textbf{ناسازگار} با هم رخ نمی‌دهند.
  \item \textbf{احتمال تجربی} $=\dfrac{N_A}{N}$؛ با تکرار زیاد به احتمال واقعی نزدیک می‌شود.
  \item سکه و تاس حافظه ندارند (خطای قمارباز).
\end{itemize}
\end{summary}

\subsection*{تمرین‌ها}

\begin{exercise}
فضای نمونهٔ آزمایش‌های زیر را بنویسید و $n(S)$ را بیابید: (الف) پرتاب یک سکه و یک تاس؛ (ب) پرتاب سه سکه؛ (ج) خانواده‌ای با دو فرزند (دختر یا پسر، به ترتیب سن).
\answer{(الف) ۱۲ \quad (ب) ۸ \quad (ج) ۴}
\end{exercise}

\begin{exercise}
در پرتاب تاس، $A$ پیشامد «عدد اول» و $B$ پیشامد «عدد کمتر از ۴» است. $A$، $B$، $A\cup B$، $A\cap B$ و $B'$ را بنویسید.
\answer{$A=\{2,3,5\}$، \ $B=\{1,2,3\}$، \ $A\cup B=\{1,2,3,5\}$، \ $A\cap B=\{2,3\}$، \ $B'=\{4,5,6\}$}
\end{exercise}

\begin{exercise}
کدام زوج‌ها ناسازگارند؟ (الف) «سکه شیر بیاید» و «سکه خط بیاید»؛ (ب) «تاس زوج بیاید» و «تاس ۶ بیاید»؛ (ج) «تاس کمتر از ۳ بیاید» و «تاس بیشتر از ۴ بیاید».
\answer{(الف) و (ج)}
\end{exercise}

\begin{exercise}
یک بسکتبالیست در ۸۰ پرتاب آزاد، ۶۰ بار موفق شده است. احتمال تجربی موفقیت او چقدر است؟ در ۲۰۰ پرتاب بعدی انتظار دارید حدوداً چند بار موفق شود؟
\answer{$0.75$؛ حدود ۱۵۰ بار}
\end{exercise}

\begin{exercise}
علی می‌گوید: «در قرعه‌کشی دیروز عدد ۷ آمد، پس امروز احتمال آمدن ۷ کمتر است.» نظر او درست است؟ توضیح دهید.
\end{exercise}

% !TEX root = ../dabirestan.tex
%% ============================================================
%%  فصل ۲ — شمارش: اصل ضرب، جایگشت و ترکیب
%% ============================================================
\chapter{شمارش: اصل ضرب، جایگشت و ترکیب}
\label{hs:ch2}

\chapterintro{\begin{itemize}
  \item با اصل جمع و اصل ضرب، تعداد حالت‌ها را بدون فهرست کردن همه‌شان بشمارید؛
  \item جایگشت (چیدن با ترتیب) و ترکیب (انتخاب بدون ترتیب) را از هم تشخیص دهید؛
  \item با فرمول‌های $n!$، $P(n,r)$ و $\binom nr$ مسئله حل کنید.
\end{itemize}}

برای حساب کردن احتمال در فصل بعد باید بتوانیم تعداد برآمدها را بشماریم. وقتی تعداد حالت‌ها کم است فهرست کردنشان آسان است، ولی مثلاً تعداد رمزهای ۴ رقمی را نمی‌توانیم یکی‌یکی بنویسیم! در این فصل ابزارهای شمردن سریع را یاد می‌گیریم.

%% ------------------------------------------------------------
\section{اصل جمع و اصل ضرب}
%% ------------------------------------------------------------

\begin{theorem}[label=hs:thm:counting]{اصل جمع و اصل ضرب}
\begin{itemize}
  \item \textbf{اصل جمع:} اگر کاری را بتوان \emph{یا} به $m$ روش \emph{یا} به $n$ روش (بدون حالت مشترک) انجام داد، روی هم $m+n$ روش وجود دارد.
  \item \textbf{اصل ضرب:} اگر کاری دو مرحله داشته باشد و مرحلهٔ اول $m$ روش و مرحلهٔ دوم (برای هر انتخاب مرحلهٔ اول) $n$ روش داشته باشد، کل کار به $m\times n$ روش انجام می‌شود. برای چند مرحله هم تعدادها را در هم ضرب می‌کنیم.
\end{itemize}
\end{theorem}

\begin{example}[label=hs:ex:clothes]{لباس پوشیدن}
سارا ۳ پیراهن و ۲ شلوار دارد. به چند طریق می‌تواند یک دست لباس (یک پیراهن و یک شلوار) بپوشد؟
\begin{solution}
مرحلهٔ اول (پیراهن) ۳ انتخاب و مرحلهٔ دوم (شلوار) ۲ انتخاب دارد؛ طبق اصل ضرب $3\times2=6$ روش. نمودار درختی همین را نشان می‌دهد:
\begin{figblock}
\begin{tikzpicture}[
  grow=left, level distance=2.8cm,
  level 1/.style={sibling distance=1.9cm},
  level 2/.style={sibling distance=0.8cm},
  every node/.style={font=\small},
  st/.style={draw=primary!60, fill=primarylight, rounded corners=2pt, inner sep=3pt},
  edge from parent/.style={draw=primary!55, thick}]
  \node[circle, fill=primary, inner sep=2pt] {}
    child { node[st] {\rl{پیراهن ۱}} child { node {\rl{شلوار ۱}} } child { node {\rl{شلوار ۲}} } }
    child { node[st] {\rl{پیراهن ۲}} child { node {\rl{شلوار ۱}} } child { node {\rl{شلوار ۲}} } }
    child { node[st] {\rl{پیراهن ۳}} child { node {\rl{شلوار ۱}} } child { node {\rl{شلوار ۲}} } };
\end{tikzpicture}

\figcaption[hs:fig:clothes]{هر مسیر یک دست لباس است: $3\times2=6$ مسیر.}
\end{figblock}
\end{solution}
\end{example}

\begin{example}[label=hs:ex:passcode]{رمز گوشی}
یک رمز ۴ رقمی از ارقام ۰ تا ۹ ساخته می‌شود. (الف) چند رمز ممکن است؟ (ب) اگر رقم تکراری مجاز نباشد چطور؟
\begin{solution}
(الف) هر رقم ۱۰ انتخاب دارد: $10\times10\times10\times10=10\,000$.\\
(ب) رقم اول ۱۰ انتخاب، دومی ۹ (چون یکی استفاده شده)، سومی ۸ و چهارمی ۷: \ $10\times9\times8\times7=5040$.
\end{solution}
\end{example}

\begin{example}[label=hs:ex:menu]{اصل جمع}
در یک رستوران ۴ نوع غذای گوشتی و ۳ نوع غذای گیاهی هست. اگر بخواهید \emph{یک} غذا انتخاب کنید، چند انتخاب دارید؟ اگر بخواهید یک غذای گوشتی \emph{و} یک غذای گیاهی بگیرید چطور؟
\begin{solution}
یک غذا (گوشتی \emph{یا} گیاهی): طبق اصل جمع $4+3=7$. \ یکی از هر نوع (گوشتی \emph{و} گیاهی): طبق اصل ضرب $4\times3=12$.
\end{solution}
\end{example}

\begin{note}
یک قاعدهٔ سرانگشتی: «\textbf{یا}» معمولاً یعنی \textbf{جمع} و «\textbf{و} سپس» معمولاً یعنی \textbf{ضرب}.
\end{note}

%% ------------------------------------------------------------
\section{جایگشت: چیدن با ترتیب}
%% ------------------------------------------------------------

\begin{definition}[label=hs:def:factorial]{فاکتوریل}
برای عدد طبیعی $n$، \term{فاکتوریل}{factorial} $n$ عبارت است از
\[ n!=n\times(n-1)\times\cdots\times2\times1, \]
و قرارداد می‌کنیم $0!=1$. مثلاً $3!=6$، \ $4!=24$ و $5!=120$.
\end{definition}

\begin{theorem}[label=hs:thm:perm]{تعداد جایگشت‌ها}
\begin{itemize}
  \item $n$ شیء متمایز را به $n!$ روش می‌توان در یک ردیف چید.
  \item تعداد روش‌های انتخاب و چیدن $r$ شیء از بین $n$ شیء (با ترتیب) برابر است با
  \[ P(n,r)=n\times(n-1)\times\cdots\times(n-r+1)=\frac{n!}{(n-r)!}. \]
\end{itemize}
\end{theorem}

\textbf{چرا؟} طبق اصل ضرب: جایگاه اول $n$ انتخاب دارد، جایگاه دوم $n-1$ انتخاب، و همین‌طور تا آخر.

\begin{example}[label=hs:ex:books]{چیدن کتاب‌ها}
به چند طریق می‌توان ۵ کتاب متفاوت را در یک قفسه کنار هم چید؟
\begin{solution}
$5!=5\times4\times3\times2\times1=120$ روش.
\end{solution}
\end{example}

\begin{example}[label=hs:ex:podium]{سکوی مسابقه}
۸ دونده در یک مسابقه شرکت کرده‌اند. به چند طریق ممکن است مقام‌های اول، دوم و سوم مشخص شوند؟
\begin{solution}
ترتیب مهم است (اول بودن با دوم بودن فرق دارد): \ $P(8,3)=8\times7\times6=336$.
\end{solution}
\end{example}

%% ------------------------------------------------------------
\section{ترکیب: انتخاب بدون ترتیب}
%% ------------------------------------------------------------

گاهی فقط مهم است \emph{چه کسانی} انتخاب شوند، نه به چه ترتیبی. مثلاً انتخاب ۲ نفر از ۴ نفر برای یک تیم: تیم «علی و مریم» همان تیم «مریم و علی» است.

\begin{definition}[label=hs:def:comb]{ترکیب}
تعداد روش‌های انتخاب $r$ شیء از بین $n$ شیء متمایز \emph{بدون در نظر گرفتن ترتیب} را با $\binom nr$ (بخوانید «$n$ انتخاب $r$») نشان می‌دهیم و آن را \term{ترکیب}{combination} می‌نامیم:
\[ \binom nr=\frac{n!}{r!\,(n-r)!}. \]
\end{definition}

\textbf{چرا؟} اگر ترتیب مهم بود، $P(n,r)$ حالت داشتیم. ولی هر گروه $r$ نفره به $r!$ ترتیب مختلف در این شمارش آمده است؛ پس بر $r!$ تقسیم می‌کنیم: \ $\binom nr=\dfrac{P(n,r)}{r!}$.

\begin{example}[label=hs:ex:team]{انتخاب تیم}
از ۴ نفر (علی، مریم، رضا، سارا) می‌خواهیم ۲ نفر را برای یک تیم انتخاب کنیم. به چند طریق؟
\begin{solution}
$\binom42=\dfrac{4!}{2!\,2!}=\dfrac{24}{4}=6$. فهرستشان: علی‌مریم، علی‌رضا، علی‌سارا، مریم‌رضا، مریم‌سارا، رضاسارا. (اگر ترتیب مهم بود، $P(4,2)=12$ حالت داشتیم؛ هر تیم دو بار شمرده می‌شد.)
\end{solution}
\end{example}

\begin{example}[label=hs:ex:handshake]{دست دادن}
در یک جلسه ۱۰ نفر حاضرند و هر دو نفر یک بار با هم دست می‌دهند. چند دست دادن انجام می‌شود؟
\begin{solution}
هر دست دادن یعنی انتخاب ۲ نفر از ۱۰ نفر (بدون ترتیب): \ $\binom{10}2=\dfrac{10\times9}{2}=45$.
\end{solution}
\end{example}

\begin{example}[label=hs:ex:committee]{کمیته با شرط}
از ۵ دختر و ۴ پسر می‌خواهیم کمیته‌ای ۳ نفره بسازیم. (الف) به چند طریق؟ (ب) اگر کمیته دقیقاً ۲ دختر داشته باشد؟ (ج) اگر \emph{دست‌کم} یک پسر داشته باشد؟
\begin{solution}
(الف) $\binom93=84$.\\
(ب) ۲ دختر از ۵ و ۱ پسر از ۴: \ $\binom52\times\binom41=10\times4=40$.\\
(ج) شمردن «دست‌کم یک پسر» مستقیماً سخت است؛ از \emph{متمم} استفاده می‌کنیم: کمیته‌هایی که \emph{هیچ} پسری ندارند (هر سه دختر) $\binom53=10$ تا هستند؛ پس جواب $84-10=74$ است.
\end{solution}
\end{example}

\begin{center}
\renewcommand{\arraystretch}{1.5}
\begin{tabular}{l|c|c}
 & \textbf{ترتیب مهم است} & \textbf{ترتیب مهم نیست}\\\hline
\textbf{نام} & جایگشت & ترکیب\\
\textbf{فرمول} & $P(n,r)=\dfrac{n!}{(n-r)!}$ & $\dbinom nr=\dfrac{n!}{r!\,(n-r)!}$\\
\textbf{مثال} & رئیس، معاون و منشی & یک کمیتهٔ ۳ نفره\\
\end{tabular}
\end{center}

\begin{warning}
قبل از هر محاسبه از خودتان بپرسید: «\emph{اگر جای دو تا از انتخاب‌ها را عوض کنم، حالت جدیدی ساخته‌ام؟}» اگر بله، ترتیب مهم است (جایگشت)؛ اگر نه، ترتیب مهم نیست (ترکیب).
\end{warning}

%% ------------------------------------------------------------
\section*{خلاصهٔ فصل ۲}
\addcontentsline{toc}{section}{خلاصه و تمرین‌ها}
%% ------------------------------------------------------------

\begin{summary}[title={\leavevmode\boxicon{accent}{$\Sigma$}\ \ خلاصهٔ فصل}]
\begin{itemize}
  \item \textbf{اصل ضرب:} مراحل پشت سر هم $\Rightarrow$ ضرب. \ \textbf{اصل جمع:} حالت‌های جدا از هم $\Rightarrow$ جمع.
  \item \textbf{جایگشت} (ترتیب مهم): $n!$ و $P(n,r)=\frac{n!}{(n-r)!}$.
  \item \textbf{ترکیب} (ترتیب مهم نیست): $\binom nr=\frac{n!}{r!(n-r)!}$.
  \item برای «دست‌کم یکی» معمولاً راحت‌تر است حالت‌های «هیچ» را از کل کم کنیم.
\end{itemize}
\end{summary}

\subsection*{تمرین‌ها}

\begin{exercise}
پلاکی از ۲ حرف (از ۳۲ حرف الفبا) و سپس ۳ رقم تشکیل شده است. چند پلاک مختلف ممکن است؟
\answer{$32\times32\times1000=1\,024\,000$}
\end{exercise}

\begin{exercise}
به چند طریق می‌توان ۶ نفر را در یک صف ایستاد؟ اگر علی بخواهد حتماً اول صف باشد چطور؟
\answer{$720$؛ \ $120$}
\end{exercise}

\begin{exercise}
از یک کلاس ۲۵ نفره، یک نماینده و یک معاون انتخاب می‌شود. به چند طریق؟
\answer{$25\times24=600$}
\end{exercise}

\begin{exercise}
در یک آزمون باید به ۶ سؤال از ۸ سؤال جواب دهید. به چند طریق می‌توانید سؤال‌ها را انتخاب کنید؟
\answer{$\binom86=28$}
\end{exercise}

\begin{exercise}
یک جعبه ۶ مهرهٔ قرمز و ۴ مهرهٔ آبی دارد. به چند طریق می‌توان ۳ مهره برداشت که (الف) همه قرمز باشند؛ (ب) ۲ قرمز و ۱ آبی باشند؛ (ج) دست‌کم یک آبی باشد؟
\answer{(الف) $20$ \quad (ب) $60$ \quad (ج) $120-20=100$}
\end{exercise}

\begin{exercise}
با حروف کلمهٔ «کتاب» (۴ حرف متفاوت) چند کلمهٔ ۴ حرفی (با معنی یا بی‌معنی) می‌توان ساخت؟ چند کلمهٔ ۲ حرفی؟
\answer{$24$؛ \ $12$}
\end{exercise}

% !TEX root = ../dabirestan.tex
%% ============================================================
%%  فصل ۳ — احتمال در فضاهای هم‌شانس و قوانین احتمال
%% ============================================================
\chapter{احتمال در فضاهای هم‌شانس و قوانین احتمال}
\label{hs:ch3}

\chapterintro{\begin{itemize}
  \item وقتی برآمدها هم‌شانس‌اند، احتمال را با فرمول $\P(A)=\frac{n(A)}{n(S)}$ حساب کنید؛
  \item با ابزارهای شمارش فصل ۲ مسئله‌های احتمال (تاس، کارت، مهره) را حل کنید؛
  \item قانون متمم و قانون جمع را بفهمید و به کار ببرید.
\end{itemize}}

%% ------------------------------------------------------------
\section{احتمال کلاسیک}
%% ------------------------------------------------------------

اگر تاسی سالم باشد، هیچ وجهی بر وجه دیگر برتری ندارد؛ می‌گوییم برآمدها \term{هم‌شانس}{equally likely}اند. در این حالت احتمال را بدون آزمایش و فقط با شمردن حساب می‌کنیم.

\begin{definition}[label=hs:def:classical]{احتمال در فضای هم‌شانس}
اگر همهٔ برآمدهای فضای نمونهٔ $S$ هم‌شانس باشند، احتمال پیشامد $A$ برابر است با
\[ \P(A)=\frac{n(A)}{n(S)}=\frac{\text{\rl{تعداد حالت‌های مطلوب}}}{\text{\rl{تعداد کل حالت‌ها}}}. \]
\end{definition}

\begin{example}[label=hs:ex:die]{تاس}
تاس سالمی پرتاب می‌شود. احتمال پیشامدهای «زوج آمدن» و «آمدن عددی بیشتر از ۴» را بیابید.
\begin{solution}
$n(S)=6$. \ زوج: $A=\{2,4,6\}$، پس $\P(A)=\tfrac36=\tfrac12$. \ بیشتر از ۴: $B=\{5,6\}$، پس $\P(B)=\tfrac26=\tfrac13$.
\end{solution}
\end{example}

\begin{example}[label=hs:ex:twocoins]{دو سکه}
دو سکهٔ سالم پرتاب می‌شوند. احتمال اینکه «یکی شیر و یکی خط» بیاید چقدر است؟
\begin{solution}
$S=\{HH,HT,TH,TT\}$ با ۴ برآمد هم‌شانس. پیشامد مطلوب $\{HT,TH\}$ دو عضو دارد؛ پس احتمال $\tfrac24=\tfrac12$ است. (اگر به اشتباه فقط سه حالت «دو شیر، دو خط، یکی از هر کدام» را در نظر می‌گرفتیم، جواب غلطِ $\tfrac13$ به دست می‌آمد؛ آن سه حالت هم‌شانس نیستند.)
\end{solution}
\end{example}

\begin{example}[label=hs:ex:twodice]{مجموع دو تاس}
دو تاس سالم (یکی سفید و یکی قرمز) پرتاب می‌شوند. احتمال اینکه مجموع دو عدد ۷ شود چقدر است؟ مجموع ۱۲ چطور؟
\begin{solution}
طبق اصل ضرب $n(S)=6\times6=36$ برآمد هم‌شانس داریم. جدول زیر مجموع هر برآمد را نشان می‌دهد:
\begin{figblock}
\begin{tikzpicture}[x=0.82cm, y=0.62cm]
  \foreach \i in {1,...,6} \foreach \j in {1,...,6} {
    \pgfmathtruncatemacro{\s}{\i+\j}
    \ifnum\s=7 \fill[accent!40] (\i-0.5,-\j+0.5) rectangle (\i+0.5,-\j-0.5); \fi
    \node[font=\small] at (\i,-\j) {\lr{\s}};
  }
  \draw[primary, thick] (0.5,-0.5) rectangle (6.5,-6.5);
  \foreach \i in {1,...,6} { \node[font=\small\bfseries, text=defcolor] at (\i,0.3) {\lr{\i}}; \node[font=\small\bfseries, text=warncolor] at (7.2,-\i) {\lr{\i}}; }
  \node[font=\scriptsize, text=defcolor] at (3.5,1.1) {\rl{تاس سفید}};
  \node[font=\scriptsize, text=warncolor] at (8.4,-3.5) {\rl{تاس قرمز}};
\end{tikzpicture}

\figcaption[hs:fig:dice-table]{جدول مجموع دو تاس؛ خانه‌های رنگی مجموع ۷ را نشان می‌دهند.}
\end{figblock}
مجموع ۷ در ۶ خانه آمده است: $\P(\fa{مجموع ۷})=\tfrac{6}{36}=\tfrac16$. \ مجموع ۱۲ فقط یک خانه دارد: $\P(\fa{مجموع ۱۲})=\tfrac1{36}$.
\end{solution}
\end{example}

\begin{warning}
فرمول $\frac{n(A)}{n(S)}$ \emph{فقط} وقتی درست است که برآمدها هم‌شانس باشند. مثلاً اگر فضای نمونهٔ مجموع دو تاس را $\{2,3,\dots,12\}$ بگیرید (۱۱ حالت)، این حالت‌ها هم‌شانس نیستند و نمی‌توانید بگویید احتمال مجموع ۷ برابر $\tfrac1{11}$ است.
\end{warning}

\subsection{مسائل با کارت و مهره}

یک دست کارت بازی ۵۲ کارت دارد: ۴ خال (\lr{$\spadesuit$} پیک، \lr{$\heartsuit$} دل، \lr{$\diamondsuit$} خشت، \lr{$\clubsuit$} گشنیز) و در هر خال ۱۳ کارت (آس، ۲ تا ۱۰، سرباز، بی‌بی، شاه).

\begin{example}[label=hs:ex:card]{یک کارت}
یک کارت به تصادف انتخاب می‌کنیم. احتمال اینکه (الف) آس باشد؛ (ب) دل باشد؛ (ج) شاهِ دل باشد چقدر است؟
\begin{solution}
$n(S)=52$. \ (الف) ۴ آس: $\tfrac4{52}=\tfrac1{13}$. \ (ب) ۱۳ دل: $\tfrac{13}{52}=\tfrac14$. \ (ج) فقط یک کارت: $\tfrac1{52}$.
\end{solution}
\end{example}

\begin{example}[label=hs:ex:urn]{دو مهره از جعبه}
جعبه‌ای ۵ مهرهٔ قرمز و ۳ مهرهٔ آبی دارد. ۲ مهره به تصادف (با هم) برمی‌داریم. احتمال اینکه (الف) هر دو قرمز باشند؛ (ب) یکی قرمز و یکی آبی باشد چقدر است؟
\begin{solution}
ترتیب مهم نیست؛ $n(S)=\binom82=28$.\\
(الف) ۲ قرمز از ۵: $\binom52=10$؛ پس $\P=\tfrac{10}{28}=\tfrac5{14}$.\\
(ب) ۱ قرمز از ۵ و ۱ آبی از ۳: $5\times3=15$؛ پس $\P=\tfrac{15}{28}$.
\end{solution}
\end{example}

\begin{example}[label=hs:ex:class]{انتخاب نماینده}
کلاسی ۱۲ دختر و ۸ پسر دارد. ۳ نفر به تصادف برای یک اردو انتخاب می‌شوند. احتمال اینکه هر سه دختر باشند چقدر است؟
\begin{solution}
$\P=\dfrac{\binom{12}3}{\binom{20}3}=\dfrac{220}{1140}=\dfrac{11}{57}\approx0.19$.
\end{solution}
\end{example}

%% ------------------------------------------------------------
\section{قوانین احتمال}
%% ------------------------------------------------------------

از فرمول $\P(A)=\frac{n(A)}{n(S)}$ چند قانون ساده ولی بسیار مهم به دست می‌آید. این قانون‌ها در حالت کلی (حتی وقتی برآمدها هم‌شانس نیستند) هم درست‌اند.

\begin{theorem}[label=hs:thm:basic]{قوانین پایهٔ احتمال}
برای هر پیشامد $A$:
\begin{enumerate}
  \item $0\le\P(A)\le1$؛
  \item $\P(S)=1$ (پیشامد حتمی) و $\P(\varnothing)=0$ (پیشامد غیرممکن)؛
  \item \textbf{قانون متمم:} \ $\P(A')=1-\P(A)$.
\end{enumerate}
\end{theorem}

\textbf{چرا قانون متمم درست است؟} هر برآمد یا در $A$ است یا در $A'$؛ پس $n(A)+n(A')=n(S)$. با تقسیم بر $n(S)$ داریم $\P(A)+\P(A')=1$.

\begin{example}[label=hs:ex:atleast]{دست‌کم یک شیر}
سه سکهٔ سالم پرتاب می‌شوند. احتمال اینکه \emph{دست‌کم یک} شیر بیاید چقدر است؟
\begin{solution}
$n(S)=2^3=8$. متمم «دست‌کم یک شیر» این است: «هیچ شیری نیاید»، یعنی فقط $TTT$. پس
\[ \P(\fa{دست‌کم یک شیر})=1-\P(TTT)=1-\tfrac18=\tfrac78. \]
\end{solution}
\end{example}

\begin{note}
هر وقت در صورت مسئله «\textbf{دست‌کم یکی}» دیدید، اول به متمم فکر کنید: «\textbf{هیچ‌کدام}» معمولاً خیلی راحت‌تر حساب می‌شود.
\end{note}

\begin{theorem}[label=hs:thm:add]{قانون جمع}
\begin{itemize}
  \item اگر $A$ و $B$ \textbf{ناسازگار} باشند: \quad $\P(A\cup B)=\P(A)+\P(B)$.
  \item در حالت کلی: \quad $\P(A\cup B)=\P(A)+\P(B)-\P(A\cap B)$.
\end{itemize}
\end{theorem}

\textbf{چرا؟} در $\P(A)+\P(B)$، برآمدهایی که هم در $A$ و هم در $B$ هستند (ناحیهٔ مشترک در نمودار ون) \emph{دو بار} شمرده شده‌اند؛ پس یک بار $\P(A\cap B)$ را کم می‌کنیم. اگر ناسازگار باشند، ناحیهٔ مشترکی نیست.

\begin{example}[label=hs:ex:acehearts]{آس یا دل}
یک کارت به تصادف انتخاب می‌کنیم. احتمال اینکه آس \emph{یا} دل باشد چقدر است؟
\begin{solution}
$\P(\fa{آس})=\tfrac4{52}$ و $\P(\fa{دل})=\tfrac{13}{52}$. ولی آسِ دل در هر دو شمرده شده است: $\P(\fa{آس دل})=\tfrac1{52}$. پس
\[ \P(\fa{آس یا دل})=\frac4{52}+\frac{13}{52}-\frac1{52}=\frac{16}{52}=\frac4{13}. \]
\end{solution}
\end{example}

\begin{example}[label=hs:ex:sports]{ورزش دانش‌آموزان}
در یک کلاس ۳۰ نفره، ۱۸ نفر فوتبال، ۱۲ نفر والیبال و ۷ نفر هر دو ورزش را بازی می‌کنند. اگر یک نفر را به تصادف انتخاب کنیم، احتمال اینکه (الف) دست‌کم یکی از این دو ورزش را بازی کند؛ (ب) هیچ‌کدام را بازی نکند چقدر است؟
\begin{solution}
(الف) $\P(F\cup V)=\tfrac{18}{30}+\tfrac{12}{30}-\tfrac{7}{30}=\tfrac{23}{30}$.
\ (ب) طبق قانون متمم: $1-\tfrac{23}{30}=\tfrac7{30}$.
نمودار ون این را نشان می‌دهد:
\begin{figblock}
\begin{tikzpicture}[scale=0.9]
  \draw[graytext, thick, rounded corners=5pt] (-2.4,-1.7) rectangle (4.2,1.7);
  \fill[setA, opacity=0.25] (0,0) circle (1.3); \fill[setB, opacity=0.25] (1.8,0) circle (1.3);
  \draw[setA!80!black, thick] (0,0) circle (1.3); \draw[setB!80!black, thick] (1.8,0) circle (1.3);
  \node[font=\small] at (-0.45,0) {\lr{11}}; \node[font=\small] at (0.9,0) {\lr{7}}; \node[font=\small] at (2.25,0) {\lr{5}};
  \node[font=\small] at (3.75,-1.35) {\lr{7}};
  \node[font=\small, text=setA!80!black] at (-0.9,1.55) {\rl{فوتبال}};
  \node[font=\small, text=setB!80!black] at (2.7,1.55) {\rl{والیبال}};
\end{tikzpicture}

\figcaption[hs:fig:sports]{۱۱ نفر فقط فوتبال، ۷ نفر هر دو، ۵ نفر فقط والیبال و ۷ نفر هیچ‌کدام.}
\end{figblock}
\end{solution}
\end{example}

\begin{moreinfo}{مسئلهٔ روز تولد}
در یک کلاس ۲۳ نفره، احتمال اینکه دست‌کم دو نفر روز تولد یکسان داشته باشند چقدر است؟ حدس بزنید! جواب شگفت‌انگیز است: \textbf{بیش از ۵۰٪}. با قانون متمم: احتمال اینکه همه روز تولد متفاوت داشته باشند
\[ \frac{365}{365}\times\frac{364}{365}\times\frac{363}{365}\times\cdots\times\frac{343}{365}\approx0.493 \]
است؛ پس احتمال دست‌کم یک تکرار حدود $1-0.493=0.507$ است. با ۵۰ نفر این احتمال به حدود ۹۷٪ می‌رسد!
\end{moreinfo}

%% ------------------------------------------------------------
\section*{خلاصهٔ فصل ۳}
\addcontentsline{toc}{section}{خلاصه و تمرین‌ها}
%% ------------------------------------------------------------

\begin{summary}[title={\leavevmode\boxicon{accent}{$\Sigma$}\ \ خلاصهٔ فصل}]
\begin{itemize}
  \item برای برآمدهای \textbf{هم‌شانس}: $\P(A)=\frac{n(A)}{n(S)}$. قبل از استفاده مطمئن شوید برآمدها واقعاً هم‌شانس‌اند.
  \item $0\le\P(A)\le1$، \ $\P(S)=1$، \ $\P(\varnothing)=0$.
  \item \textbf{قانون متمم:} $\P(A')=1-\P(A)$؛ بهترین ابزار برای «دست‌کم یکی».
  \item \textbf{قانون جمع:} $\P(A\cup B)=\P(A)+\P(B)-\P(A\cap B)$؛ برای پیشامدهای ناسازگار جملهٔ آخر صفر است.
\end{itemize}
\end{summary}

\subsection*{تمرین‌ها}

\begin{exercise}
دو تاس سالم پرتاب می‌شوند. احتمال اینکه (الف) مجموع ۸ شود؛ (ب) دو عدد برابر بیایند (جفت)؛ (ج) مجموع کمتر از ۵ شود چقدر است؟
\answer{(الف) $\tfrac5{36}$ \quad (ب) $\tfrac16$ \quad (ج) $\tfrac16$}
\end{exercise}

\begin{exercise}
یک کارت از دست ۵۲ کارتی انتخاب می‌شود. احتمال اینکه (الف) کارت تصویری (سرباز، بی‌بی یا شاه) باشد؛ (ب) قرمز (دل یا خشت) یا شاه باشد چقدر است؟
\answer{(الف) $\tfrac{12}{52}=\tfrac3{13}$ \quad (ب) $\tfrac{26+4-2}{52}=\tfrac7{13}$}
\end{exercise}

\begin{exercise}
از اعداد ۱ تا ۳۰ یک عدد به تصادف انتخاب می‌شود. احتمال اینکه بر ۲ \emph{یا} بر ۳ بخش‌پذیر باشد چقدر است؟
\answer{$\tfrac{15+10-5}{30}=\tfrac23$}
\end{exercise}

\begin{exercise}
جعبه‌ای ۴ مهرهٔ سفید و ۶ مهرهٔ سیاه دارد. ۳ مهره به تصادف برمی‌داریم. احتمال اینکه (الف) هر سه سیاه باشند؛ (ب) دست‌کم یک سفید باشد چقدر است؟
\answer{(الف) $\tfrac{20}{120}=\tfrac16$ \quad (ب) $\tfrac56$}
\end{exercise}

\begin{exercise}
اگر $\P(A)=0.5$، $\P(B)=0.4$ و $\P(A\cap B)=0.2$، مقادیر $\P(A\cup B)$ و $\P\bigl((A\cup B)'\bigr)$ را بیابید.
\answer{$0.7$ و $0.3$}
\end{exercise}

\begin{exercise}
چهار سکهٔ سالم پرتاب می‌شوند. احتمال اینکه دقیقاً ۲ شیر بیاید چقدر است؟ دست‌کم یک خط چطور؟
\answer{$\tfrac{\binom42}{16}=\tfrac6{16}=\tfrac38$؛ \ $1-\tfrac1{16}=\tfrac{15}{16}$}
\end{exercise}

% !TEX root = ../dabirestan.tex
%% ============================================================
%%  فصل ۴ — احتمال شرطی و نمودار درختی
%% ============================================================
\chapter{احتمال شرطی و نمودار درختی}
\label{hs:ch4}

\chapterintro{\begin{itemize}
  \item بفهمید «اطلاعات جدید» چطور احتمال را تغییر می‌دهد؛
  \item احتمال شرطی را با فرمول و با «کوچک کردن فضای نمونه» حساب کنید؛
  \item با قانون ضرب و نمودار درختی، احتمال آزمایش‌های چندمرحله‌ای را بیابید.
\end{itemize}}

%% ------------------------------------------------------------
\section{وقتی چیزی می‌دانیم}
%% ------------------------------------------------------------

\begin{example}[label=hs:ex:hint]{یک راهنمایی کوچک}
دوستتان تاسی می‌اندازد و نتیجه را پنهان می‌کند. احتمال اینکه ۶ آمده باشد $\tfrac16$ است. حالا به شما می‌گوید: «عدد زوج آمده است.» حالا احتمال ۶ چقدر است؟
\begin{solution}
با این راهنمایی، برآمدهای فرد کنار می‌روند و فقط $\{2,4,6\}$ باقی می‌ماند. این سه حالت هنوز هم‌شانس‌اند و یکی از آن‌ها ۶ است؛ پس احتمال حالا $\tfrac13$ است. \emph{اطلاعات جدید، احتمال را عوض کرد.}
\end{solution}
\end{example}

این احتمالِ تازه را «احتمال ۶ \emph{به شرط} زوج بودن» می‌نامیم و می‌نویسیم $\P(\fa{۶}\mid\fa{زوج})$. خط عمودی $\mid$ را «به شرط» می‌خوانیم.

\begin{figblock}
\begin{tikzpicture}[scale=0.85]
  \draw[graytext, thick, rounded corners=5pt] (-2,-1.5) rectangle (3.2,1.5);
  \fill[setA, opacity=0.25] (0,0) circle (1.1); \fill[setB, opacity=0.3] (1.3,0) circle (1.1);
  \draw[setA!80!black, thick] (0,0) circle (1.1); \draw[setB!80!black, thick] (1.3,0) circle (1.1);
  \node at (-0.45,0.2) {\lr{$A$}}; \node at (1.85,0.2) {\lr{$B$}}; \node at (2.9,1.2) {\lr{$S$}};
  \node[font=\small] at (0.6,-1.95) {\rl{قبل از اطلاع: کل $S$}};
  \draw[-{Stealth[length=3mm]}, line width=1.1pt, accent] (3.6,0) -- node[above, font=\small]{\rl{$B$ رخ داد}} (5.4,0);
  \begin{scope}[xshift=7.4cm]
    \draw[graytext!40, thick, dashed, rounded corners=5pt] (-2,-1.5) rectangle (3.2,1.5);
    \begin{scope}\clip (1.3,0) circle (1.1); \fill[setA!70] (0,0) circle (1.1); \end{scope}
    \draw[setA!40, thick, dashed] (0,0) circle (1.1);
    \draw[setB!80!black, line width=1.6pt] (1.3,0) circle (1.1);
    \node[font=\scriptsize] at (0.62,0) {\lr{$A\!\cap\!B$}}; \node at (2.05,0.35) {\lr{$B$}};
    \node[font=\small] at (0.6,-1.95) {\rl{بعد از اطلاع: $B$ فضای نمونهٔ جدید است}};
  \end{scope}
\end{tikzpicture}

\figcaption[hs:fig:cond]{احتمال شرطی یعنی فضای نمونه را به $B$ کوچک کنیم و ببینیم چه سهمی از $B$ در $A$ هم هست.}
\end{figblock}

%% ------------------------------------------------------------
\section{تعریف احتمال شرطی}
%% ------------------------------------------------------------

\begin{definition}[label=hs:def:cond]{احتمال شرطی}
اگر $\P(B)>0$، احتمال $A$ به شرط $B$ برابر است با
\[ \P(A\mid B)=\frac{\P(A\cap B)}{\P(B)}. \]
اگر برآمدها هم‌شانس باشند، این همان است که بگوییم \ $\P(A\mid B)=\dfrac{n(A\cap B)}{n(B)}$.
\end{definition}

\begin{example}[label=hs:ex:glasses]{جدول دانش‌آموزان}
در کلاسی ۳۰ نفره وضعیت عینکی بودن دانش‌آموزان این است:
\begin{center}
\renewcommand{\arraystretch}{1.35}
\begin{tabular}{l|cc|c}
 & عینکی & بدون عینک & جمع\\\hline
دختر & $6$ & $10$ & $16$\\
پسر & $4$ & $10$ & $14$\\\hline
جمع & $10$ & $20$ & $30$
\end{tabular}
\end{center}
یک نفر به تصادف انتخاب می‌شود. احتمال‌های زیر را بیابید: (الف) عینکی باشد؛ (ب) عینکی باشد به شرط اینکه دختر باشد؛ (ج) دختر باشد به شرط اینکه عینکی باشد.
\begin{solution}
(الف) $\P(\fa{عینکی})=\tfrac{10}{30}=\tfrac13$.\\
(ب) فقط ۱۶ دختر را در نظر می‌گیریم که ۶ نفرشان عینکی‌اند: $\P(\fa{عینکی}\mid\fa{دختر})=\tfrac{6}{16}=\tfrac38$.\\
(ج) فقط ۱۰ عینکی را در نظر می‌گیریم که ۶ نفرشان دخترند: $\P(\fa{دختر}\mid\fa{عینکی})=\tfrac{6}{10}=\tfrac35$.
\end{solution}
\end{example}

\begin{warning}
$\P(A\mid B)$ با $\P(B\mid A)$ فرق دارد! در مثال بالا $\P(\fa{عینکی}\mid\fa{دختر})=\tfrac38$ ولی $\P(\fa{دختر}\mid\fa{عینکی})=\tfrac35$. همیشه دقت کنید \emph{کدام پیشامد شرط است} (بعد از خط عمودی).
\end{warning}

\begin{example}[label=hs:ex:dice-cond]{دو تاس با شرط}
دو تاس سالم پرتاب می‌شوند و می‌دانیم مجموع آن‌ها ۸ شده است. احتمال اینکه یکی از تاس‌ها ۶ باشد چقدر است؟
\begin{solution}
برآمدهای با مجموع ۸: $\{(2,6),(3,5),(4,4),(5,3),(6,2)\}$؛ یعنی $n(B)=5$. از این‌ها دو تا ۶ دارند. پس احتمال $\tfrac25$ است. (بدون اطلاع، احتمال اینکه دست‌کم یک تاس ۶ باشد $\tfrac{11}{36}\approx0.31$ بود؛ حالا $0.4$ شد.)
\end{solution}
\end{example}

%% ------------------------------------------------------------
\section{قانون ضرب}
%% ------------------------------------------------------------

اگر فرمول احتمال شرطی را در $\P(B)$ ضرب کنیم، قانونی به دست می‌آید که برای آزمایش‌های چندمرحله‌ای بسیار مفید است:

\begin{theorem}[label=hs:thm:mult]{قانون ضرب}
\[ \P(A\cap B)=\P(A)\times\P(B\mid A). \]
یعنی: احتمال اینکه «اول $A$ و بعد $B$» رخ دهد، برابر است با احتمال $A$ ضرب در احتمال $B$ وقتی می‌دانیم $A$ رخ داده.
\end{theorem}

\begin{example}[label=hs:ex:two-reds]{دو مهره پشت سر هم}
جعبه‌ای ۵ مهرهٔ قرمز و ۳ مهرهٔ آبی دارد. دو مهره را \emph{یکی پس از دیگری} و \emph{بدون برگرداندن} برمی‌داریم. احتمال اینکه هر دو قرمز باشند چقدر است؟
\begin{solution}
احتمال قرمز بودن اولی $\tfrac58$ است. اگر اولی قرمز باشد، ۴ قرمز از ۷ مهره باقی می‌ماند؛ پس احتمال قرمز بودن دومی به شرط قرمز بودن اولی $\tfrac47$ است. طبق قانون ضرب:
\[ \P(\fa{هر دو قرمز})=\frac58\times\frac47=\frac{20}{56}=\frac5{14}. \]
(این همان جواب مثال \ref{hs:ex:urn} در فصل ۳ است که با ترکیب حساب کردیم.)
\end{solution}
\end{example}

\begin{example}[label=hs:ex:aces]{دو آس}
از یک دست ۵۲ کارتی، دو کارت پشت سر هم (بدون برگرداندن) برمی‌داریم. احتمال اینکه هر دو آس باشند چقدر است؟
\begin{solution}
$\P=\dfrac{4}{52}\times\dfrac{3}{51}=\dfrac{12}{2652}=\dfrac{1}{221}$.
\end{solution}
\end{example}

%% ------------------------------------------------------------
\section{نمودار درختی}
%% ------------------------------------------------------------

نمودار درختی بهترین ابزار برای دیدن همهٔ حالت‌های یک آزمایش چندمرحله‌ای است. روی هر شاخه احتمالِ آن مرحله (به شرط مراحل قبل) را می‌نویسیم.

\begin{figblock}
\begin{tikzpicture}[
  grow=left, level distance=3.1cm,
  level 1/.style={sibling distance=2.5cm},
  level 2/.style={sibling distance=1.15cm},
  every node/.style={font=\small},
  root/.style={circle, fill=primary, inner sep=2pt},
  st/.style={draw=primary!60, fill=primarylight, rounded corners=2pt, inner sep=3pt},
  edge from parent/.style={draw=primary!55, thick},
  lab/.style={font=\scriptsize, fill=white, inner sep=1pt}]
  \node[root] {}
    child { node[st] {\rl{قرمز}}
      child { node[st] (rr) {\rl{قرمز}} edge from parent node[lab, above] {\lr{$\frac{4}{7}$}} }
      child { node[st] (rb) {\rl{آبی}} edge from parent node[lab, below] {\lr{$\frac{3}{7}$}} }
      edge from parent node[lab, above] {\lr{$\frac{5}{8}$}} }
    child { node[st] {\rl{آبی}}
      child { node[st] (br) {\rl{قرمز}} edge from parent node[lab, above] {\lr{$\frac{5}{7}$}} }
      child { node[st] (bb) {\rl{آبی}} edge from parent node[lab, below] {\lr{$\frac{2}{7}$}} }
      edge from parent node[lab, below] {\lr{$\frac{3}{8}$}} };
  \node[anchor=east, font=\small, text=accent!80!black] at ([xshift=-4mm]rr.west) {\lr{$\frac58\cdot\frac47=\frac{20}{56}$}};
  \node[anchor=east, font=\small] at ([xshift=-4mm]rb.west) {\lr{$\frac58\cdot\frac37=\frac{15}{56}$}};
  \node[anchor=east, font=\small] at ([xshift=-4mm]br.west) {\lr{$\frac38\cdot\frac57=\frac{15}{56}$}};
  \node[anchor=east, font=\small] at ([xshift=-4mm]bb.west) {\lr{$\frac38\cdot\frac27=\frac{6}{56}$}};
  \node[font=\scriptsize, text=graytext] at (-3.1,2.1) {\rl{مهرهٔ اول}};
  \node[font=\scriptsize, text=graytext] at (-6.2,2.1) {\rl{مهرهٔ دوم}};
\end{tikzpicture}

\figcaption[hs:fig:tree-urn]{نمودار درختی برداشتن دو مهره از جعبه‌ای با ۵ قرمز و ۳ آبی (بدون برگرداندن).}
\end{figblock}

\begin{theorem}[label=hs:thm:tree]{قاعده‌های نمودار درختی}
\begin{enumerate}
  \item احتمال شاخه‌هایی که از یک نقطه جدا می‌شوند، جمعشان ۱ است.
  \item احتمال هر مسیر = \textbf{ضرب} احتمال‌های روی شاخه‌های آن (قانون ضرب).
  \item احتمال یک پیشامد = \textbf{جمع} احتمال همهٔ مسیرهایی که به آن پیشامد می‌رسند.
\end{enumerate}
\end{theorem}

\begin{example}[label=hs:ex:same-color]{هم‌رنگ}
در شکل \ref{hs:fig:tree-urn}، احتمال اینکه دو مهره هم‌رنگ باشند چقدر است؟ احتمال اینکه مهرهٔ دوم قرمز باشد چطور؟
\begin{solution}
هم‌رنگ: دو مسیر «قرمز-قرمز» و «آبی-آبی»: \ $\tfrac{20}{56}+\tfrac{6}{56}=\tfrac{26}{56}=\tfrac{13}{28}$.\\
مهرهٔ دوم قرمز: دو مسیر «قرمز-قرمز» و «آبی-قرمز»: \ $\tfrac{20}{56}+\tfrac{15}{56}=\tfrac{35}{56}=\tfrac58$. جالب است که همان احتمال قرمز بودن مهرهٔ اول است! (اگر ندانیم اولی چه بوده، دومی هیچ فرقی با اولی ندارد.)
\end{solution}
\end{example}

\begin{example}[label=hs:ex:exam]{قبولی در دو مرحله}
در یک مسابقه، ۴۰٪ شرکت‌کنندگان از مرحلهٔ اول قبول می‌شوند و از قبول‌شدگان مرحلهٔ اول، ۲۵٪ در مرحلهٔ دوم هم قبول می‌شوند. احتمال اینکه یک شرکت‌کننده هر دو مرحله را قبول شود چقدر است؟
\begin{solution}
$\P(\fa{قبول اول و دوم})=\P(\fa{قبول اول})\times\P(\fa{قبول دوم}\mid\fa{قبول اول})=0.4\times0.25=0.1$؛ یعنی ۱۰٪.
\end{solution}
\end{example}

%% ------------------------------------------------------------
\section*{خلاصهٔ فصل ۴}
\addcontentsline{toc}{section}{خلاصه و تمرین‌ها}
%% ------------------------------------------------------------

\begin{summary}[title={\leavevmode\boxicon{accent}{$\Sigma$}\ \ خلاصهٔ فصل}]
\begin{itemize}
  \item \textbf{احتمال شرطی:} $\P(A\mid B)=\frac{\P(A\cap B)}{\P(B)}$؛ یعنی $B$ فضای نمونهٔ جدید می‌شود. در حالت هم‌شانس: $\frac{n(A\cap B)}{n(B)}$.
  \item $\P(A\mid B)$ و $\P(B\mid A)$ معمولاً برابر نیستند.
  \item \textbf{قانون ضرب:} $\P(A\cap B)=\P(A)\,\P(B\mid A)$.
  \item \textbf{نمودار درختی:} روی شاخه‌ها ضرب، بین مسیرها جمع.
\end{itemize}
\end{summary}

\subsection*{تمرین‌ها}

\begin{exercise}
تاسی سالم پرتاب می‌شود. احتمال اینکه عدد ۲ آمده باشد، به شرط اینکه بدانیم عدد کمتر از ۴ است، چقدر است؟ به شرط اینکه بدانیم عدد اول است چطور؟
\answer{$\tfrac13$؛ \ $\tfrac13$}
\end{exercise}

\begin{exercise}
در مثال \ref{hs:ex:glasses}، احتمال اینکه دانش‌آموز انتخاب‌شده پسر باشد به شرط اینکه عینکی نباشد چقدر است؟
\answer{$\tfrac{10}{20}=\tfrac12$}
\end{exercise}

\begin{exercise}
دو کارت پشت سر هم و بدون برگرداندن از دست ۵۲ کارتی برمی‌داریم. احتمال اینکه هر دو دل باشند چقدر است؟
\answer{$\tfrac{13}{52}\times\tfrac{12}{51}=\tfrac1{17}$}
\end{exercise}

\begin{exercise}
جعبه‌ای ۳ مهرهٔ سفید و ۲ مهرهٔ سیاه دارد. دو مهره پشت سر هم و بدون برگرداندن برمی‌داریم. با رسم نمودار درختی احتمال‌های (الف) هر دو سفید؛ (ب) یکی سفید و یکی سیاه را بیابید.
\answer{(الف) $\tfrac3{10}$ \quad (ب) $\tfrac35$}
\end{exercise}

\begin{exercise}
اگر $\P(A)=0.6$، $\P(B)=0.5$ و $\P(A\cap B)=0.3$، مقادیر $\P(A\mid B)$ و $\P(B\mid A)$ را بیابید.
\answer{$0.6$ و $0.5$}
\end{exercise}

\begin{exercise}
۷۰٪ دانش‌آموزان یک مدرسه در کلاس تقویتی شرکت می‌کنند. از شرکت‌کنندگان در کلاس تقویتی ۸۰٪ و از بقیه ۵۰٪ در امتحان نمرهٔ بالای ۱۵ می‌گیرند. احتمال اینکه دانش‌آموزی هم در کلاس شرکت کرده باشد و هم نمرهٔ بالای ۱۵ بگیرد چقدر است؟
\answer{$0.7\times0.8=0.56$}
\end{exercise}

% !TEX root = ../dabirestan.tex
%% ============================================================
%%  فصل ۵ — پیشامدهای مستقل
%% ============================================================
\chapter{پیشامدهای مستقل}
\label{hs:ch5}

\chapterintro{\begin{itemize}
  \item بفهمید دو پیشامد «مستقل» یعنی چه و آن را آزمایش کنید؛
  \item «مستقل» را با «ناسازگار» اشتباه نگیرید؛
  \item احتمال آزمایش‌های تکراری (مثل چند پرتاب سکه یا چند پرتاب آزاد) را حساب کنید.
\end{itemize}}

%% ------------------------------------------------------------
\section{وقتی یک پیشامد چیزی دربارهٔ دیگری نمی‌گوید}
%% ------------------------------------------------------------

اگر یک سکه بیندازید و یک تاس، نتیجهٔ سکه هیچ تأثیری روی تاس ندارد. دانستن اینکه سکه شیر آمده، احتمال ۶ آمدن تاس را عوض نمی‌کند. چنین پیشامدهایی را \emph{مستقل} می‌نامیم.

\begin{definition}[label=hs:def:indep]{پیشامدهای مستقل}
دو پیشامد $A$ و $B$ را \term{مستقل}{independent} می‌نامیم اگر
\[ \P(A\cap B)=\P(A)\times\P(B). \]
اگر $\P(B)>0$، این همان است که بگوییم $\P(A\mid B)=\P(A)$؛ یعنی دانستن $B$ احتمال $A$ را عوض نمی‌کند. دو پیشامدی که مستقل نباشند \emph{وابسته} نام دارند.
\end{definition}

\begin{example}[label=hs:ex:coin-die]{سکه و تاس}
یک سکه و یک تاس می‌اندازیم. احتمال اینکه سکه شیر و تاس ۶ بیاید چقدر است؟
\begin{solution}
فضای نمونه $2\times6=12$ برآمد هم‌شانس دارد و فقط یکی از آن‌ها ($H$ و ۶) مطلوب است؛ پس احتمال $\tfrac1{12}$ است. می‌بینیم که $\tfrac1{12}=\tfrac12\times\tfrac16$؛ پس این دو پیشامد مستقل‌اند، همان‌طور که انتظار داشتیم.
\end{solution}
\end{example}

\begin{example}[label=hs:ex:die-indep]{آیا مستقل‌اند؟}
تاسی سالم پرتاب می‌شود. $A=\{2,4,6\}$ (زوج)، $B=\{1,2,3,4\}$ و $C=\{1,2,3\}$. آیا $A$ و $B$ مستقل‌اند؟ $A$ و $C$ چطور؟
\begin{solution}
$\P(A)=\tfrac12$، $\P(B)=\tfrac23$ و $A\cap B=\{2,4\}$ با احتمال $\tfrac13$. چون $\tfrac12\times\tfrac23=\tfrac13$، \textbf{$A$ و $B$ مستقل‌اند}: نیمی از اعضای $B$ زوج‌اند، همان‌طور که نیمی از کل اعداد زوج‌اند.\\
اما $A\cap C=\{2\}$ با احتمال $\tfrac16$ و $\P(A)\P(C)=\tfrac12\times\tfrac12=\tfrac14$. چون برابر نیستند، \textbf{$A$ و $C$ وابسته‌اند}: در $\{1,2,3\}$ فقط یک‌سوم اعداد زوج‌اند.
\end{solution}
\end{example}

\begin{example}[label=hs:ex:replace]{با برگرداندن و بدون برگرداندن}
جعبه‌ای ۵ مهرهٔ قرمز و ۳ مهرهٔ آبی دارد. دو مهره پشت سر هم برمی‌داریم. احتمال اینکه هر دو قرمز باشند را در دو حالت حساب کنید: (الف) مهرهٔ اول را \emph{برمی‌گردانیم}؛ (ب) \emph{برنمی‌گردانیم}.
\begin{solution}
(الف) با برگرداندن، جعبه دوباره مثل اول می‌شود و برداشت دوم مستقل از اولی است: $\tfrac58\times\tfrac58=\tfrac{25}{64}\approx0.39$.\\
(ب) بدون برگرداندن، برداشت دوم به اولی \emph{وابسته} است: $\tfrac58\times\tfrac47=\tfrac{5}{14}\approx0.36$.
\end{solution}
\end{example}

\begin{warning}
\textbf{مستقل با ناسازگار فرق دارد.} اگر $A$ و $B$ ناسازگار باشند (با هم رخ ندهند) و هر دو احتمال مثبت داشته باشند، آن‌ها \emph{وابسته‌اند}: اگر بدانیم $B$ رخ داده، مطمئن می‌شویم $A$ رخ نداده است! مثلاً در پرتاب تاس، «زوج آمدن» و «فرد آمدن» ناسازگارند ولی کاملاً به هم وابسته‌اند.
\end{warning}

%% ------------------------------------------------------------
\section{آزمایش‌های تکراری}
%% ------------------------------------------------------------

وقتی یک آزمایش را چند بار تکرار می‌کنیم و نتیجه‌ها روی هم اثر ندارند (مثل پرتاب‌های پیاپی سکه)، احتمال هر دنباله از نتیجه‌ها \emph{حاصل‌ضرب} احتمال‌هاست. این قانون برای بیش از دو پیشامد مستقل هم برقرار است:
\[ \P(A_1\cap A_2\cap\cdots\cap A_n)=\P(A_1)\times\P(A_2)\times\cdots\times\P(A_n). \]

\begin{example}[label=hs:ex:freethrow]{پرتاب آزاد}
بسکتبالیستی هر پرتاب آزاد را با احتمال $0.7$ گل می‌کند و پرتاب‌هایش مستقل از هم‌اند. اگر سه پرتاب انجام دهد، احتمال اینکه (الف) هر سه گل شود؛ (ب) هیچ‌کدام گل نشود؛ (ج) دست‌کم یکی گل شود؛ (د) دقیقاً دو تا گل شود چقدر است؟
\begin{solution}
(الف) $0.7\times0.7\times0.7=0.343$.\\
(ب) احتمال خطا در هر پرتاب $0.3$ است: $0.3^3=0.027$.\\
(ج) طبق قانون متمم: $1-0.027=0.973$.\\
(د) دنباله‌های با دقیقاً دو گل: «گل، گل، خطا»، «گل، خطا، گل» و «خطا، گل، گل». هر کدام احتمال $0.7\times0.7\times0.3=0.147$ دارد؛ پس جواب $3\times0.147=0.441$ است.
\end{solution}
\end{example}

\begin{example}[label=hs:ex:demere]{معمای شوالیه دومِره}
یک قمارباز فرانسوی در قرن هفدهم می‌پرسید: کدام محتمل‌تر است؟ (الف) در ۴ پرتاب یک تاس دست‌کم یک ۶ بیاید؛ (ب) در ۲۴ پرتاب دو تاس دست‌کم یک جفت‌شش بیاید.
\begin{solution}
(الف) احتمال «هیچ ۶ نیاید» در ۴ پرتاب مستقل برابر $(\tfrac56)^4$ است؛ پس احتمال مطلوب $1-(\tfrac56)^4\approx0.518$.\\
(ب) در هر پرتاب دو تاس، احتمال جفت‌شش $\tfrac1{36}$ است؛ پس احتمال مطلوب $1-(\tfrac{35}{36})^{24}\approx0.491$.\\
پس (الف) کمی محتمل‌تر است. پاسکال با همین محاسبه به این سؤال جواب داد؛ و این نامه‌نگاری‌ها آغاز علم احتمال شد!
\end{solution}
\end{example}

\begin{example}[label=hs:ex:lights]{چراغ‌های راهنمایی}
در مسیر مدرسه دو چراغ راهنمایی هست که مستقل از هم کار می‌کنند. احتمال سبز بودن چراغ اول $0.6$ و چراغ دوم $0.5$ است. احتمال اینکه (الف) هر دو سبز باشند؛ (ب) دست‌کم یکی قرمز باشد چقدر است؟
\begin{solution}
(الف) $0.6\times0.5=0.3$. \quad (ب) متمم «هر دو سبز» است: $1-0.3=0.7$.
\end{solution}
\end{example}

\begin{moreinfo}{سیستم‌های سری و موازی}
مهندسان برای بالا بردن اطمینان یک سیستم از قطعات «موازی» (یدکی) استفاده می‌کنند. فرض کنید هر قطعه مستقل از بقیه با احتمال $0.9$ سالم بماند.
\begin{itemize}
  \item \textbf{سه قطعهٔ سری} (همه باید سالم باشند، مثل حلقه‌های یک زنجیر): $0.9^3=0.729$.
  \item \textbf{سه قطعهٔ موازی} (یکی سالم بماند کافی است): $1-0.1^3=0.999$.
\end{itemize}
به همین دلیل هواپیماها چند موتور و چند سیستم کنترل موازی دارند.
\end{moreinfo}

\begin{warning}
یادآوری \textbf{خطای قمارباز}: چون پرتاب‌های سکه مستقل‌اند، احتمال شیر در پرتاب ششم همیشه $\tfrac12$ است؛ حتی اگر پنج پرتاب قبلی همه خط آمده باشند.
\end{warning}

%% ------------------------------------------------------------
\section*{خلاصهٔ فصل ۵}
\addcontentsline{toc}{section}{خلاصه و تمرین‌ها}
%% ------------------------------------------------------------

\begin{summary}[title={\leavevmode\boxicon{accent}{$\Sigma$}\ \ خلاصهٔ فصل}]
\begin{itemize}
  \item \textbf{مستقل:} $\P(A\cap B)=\P(A)\P(B)$، یا معادلاً $\P(A\mid B)=\P(A)$.
  \item ناسازگار $\ne$ مستقل؛ پیشامدهای ناسازگار (با احتمال مثبت) وابسته‌اند.
  \item برداشت \emph{با برگرداندن} مستقل است و \emph{بدون برگرداندن} وابسته.
  \item در آزمایش‌های تکراری مستقل، احتمال‌ها در هم ضرب می‌شوند؛ «دست‌کم یکی» $=1-$ «هیچ‌کدام».
\end{itemize}
\end{summary}

\subsection*{تمرین‌ها}

\begin{exercise}
سکه‌ای سالم ۴ بار پرتاب می‌شود. احتمال اینکه (الف) همه شیر بیاید؛ (ب) دست‌کم یک خط بیاید چقدر است؟
\answer{(الف) $\tfrac1{16}$ \quad (ب) $\tfrac{15}{16}$}
\end{exercise}

\begin{exercise}
دو تاس پرتاب می‌شوند. $A$: «تاس اول زوج است» و $B$: «مجموع دو تاس ۷ است». آیا $A$ و $B$ مستقل‌اند؟
\answer{بله؛ $\P(A\cap B)=\tfrac3{36}=\tfrac12\times\tfrac16$}
\end{exercise}

\begin{exercise}
اگر $A$ و $B$ مستقل باشند و $\P(A)=0.4$ و $\P(B)=0.5$، مقادیر $\P(A\cap B)$ و $\P(A\cup B)$ را بیابید.
\answer{$0.2$ و $0.7$}
\end{exercise}

\begin{exercise}
احتمال اینکه علی به یک سؤال تستی چهارگزینه‌ای به صورت شانسی درست جواب دهد $\tfrac14$ است. اگر به ۳ سؤال شانسی جواب دهد، احتمال اینکه (الف) همه را درست بزند؛ (ب) دست‌کم یکی را درست بزند چقدر است؟
\answer{(الف) $\tfrac1{64}$ \quad (ب) $1-\tfrac{27}{64}=\tfrac{37}{64}$}
\end{exercise}

\begin{exercise}
احتمال باران در هر یک از روزهای شنبه، یکشنبه و دوشنبه $0.2$ است و روزها مستقل فرض می‌شوند. احتمال اینکه در این سه روز دست‌کم یک روز باران بیاید چقدر است؟
\answer{$1-0.8^3=0.488$}
\end{exercise}

% !TEX root = ../dabirestan.tex
%% ============================================================
%%  فصل ۶ — قانون احتمال کل و قضیهٔ بیز
%% ============================================================
\chapter{قانون احتمال کل و قضیهٔ بیز}
\label{hs:ch6}

\chapterintro{\begin{itemize}
  \item با «شکستن مسئله به حالت‌ها» احتمال یک پیشامد را حساب کنید (قانون احتمال کل)؛
  \item با قضیهٔ بیز از «نتیجه» به «علت» برگردید؛
  \item بفهمید چرا یک آزمایش پزشکیِ «۹۵٪ دقیق» ممکن است بیشتر وقت‌ها اشتباه کند!
\end{itemize}}

%% ------------------------------------------------------------
\section{قانون احتمال کل}
%% ------------------------------------------------------------

گاهی احتمال یک پیشامد به این بستگی دارد که «کدام حالت» پیش آمده است. ایدهٔ قانون احتمال کل ساده است: \key{مسئله را به حالت‌های جدا از هم بشکن، هر حالت را جدا حساب کن و نتیجه‌ها را جمع بزن.}

\begin{example}[label=hs:ex:two-boxes]{دو جعبه}
جعبهٔ ۱ سه مهرهٔ قرمز و دو مهرهٔ آبی دارد و جعبهٔ ۲ یک مهرهٔ قرمز و چهار مهرهٔ آبی. سکه‌ای می‌اندازیم: اگر شیر آمد از جعبهٔ ۱ و اگر خط آمد از جعبهٔ ۲ یک مهره برمی‌داریم. احتمال اینکه مهره قرمز باشد چقدر است؟
\begin{solution}
نمودار درختی را رسم می‌کنیم:
\begin{figblock}
\begin{tikzpicture}[
  grow=left, level distance=3.2cm,
  level 1/.style={sibling distance=2.5cm},
  level 2/.style={sibling distance=1.15cm},
  every node/.style={font=\small},
  root/.style={circle, fill=primary, inner sep=2pt},
  st/.style={draw=primary!60, fill=primarylight, rounded corners=2pt, inner sep=3pt},
  red/.style={draw=accent, fill=notebg, rounded corners=2pt, inner sep=3pt},
  edge from parent/.style={draw=primary!55, thick},
  lab/.style={font=\scriptsize, fill=white, inner sep=1pt}]
  \node[root] {}
    child { node[st] {\rl{جعبهٔ ۱}}
      child { node[red] (r1) {\rl{قرمز}} edge from parent node[lab, above] {\lr{$\frac35$}} }
      child { node[st] {\rl{آبی}} edge from parent node[lab, below] {\lr{$\frac25$}} }
      edge from parent node[lab, above] {\lr{$\frac12$}} }
    child { node[st] {\rl{جعبهٔ ۲}}
      child { node[red] (r2) {\rl{قرمز}} edge from parent node[lab, above] {\lr{$\frac15$}} }
      child { node[st] {\rl{آبی}} edge from parent node[lab, below] {\lr{$\frac45$}} }
      edge from parent node[lab, below] {\lr{$\frac12$}} };
  \node[anchor=east, font=\small, text=accent!80!black] at ([xshift=-4mm]r1.west) {\lr{$\frac12\cdot\frac35=\frac3{10}$}};
  \node[anchor=east, font=\small, text=accent!80!black] at ([xshift=-4mm]r2.west) {\lr{$\frac12\cdot\frac15=\frac1{10}$}};
\end{tikzpicture}

\figcaption[hs:fig:two-boxes]{دو مسیر به «قرمز» می‌رسند؛ احتمال قرمز = جمع احتمال این دو مسیر.}
\end{figblock}
\[ \P(\fa{قرمز})=\frac12\cdot\frac35+\frac12\cdot\frac15=\frac3{10}+\frac1{10}=\frac25. \]
\end{solution}
\end{example}

\begin{definition}[label=hs:def:partition]{افراز}
پیشامدهای $B_1,B_2,\dots,B_k$ یک \term{افراز}{partition} از فضای نمونه‌اند اگر دوبه‌دو ناسازگار باشند و اجتماعشان کل $S$ باشد؛ یعنی در هر آزمایش \emph{دقیقاً یکی} از آن‌ها رخ دهد. (مثال بالا: «جعبهٔ ۱» و «جعبهٔ ۲».)
\end{definition}

\begin{theorem}[label=hs:thm:total]{قانون احتمال کل}
اگر $B_1,\dots,B_k$ افرازی از $S$ باشند، برای هر پیشامد $A$:
\[ \P(A)=\P(B_1)\,\P(A\mid B_1)+\P(B_2)\,\P(A\mid B_2)+\cdots+\P(B_k)\,\P(A\mid B_k). \]
\end{theorem}

\textbf{چرا؟} پیشامد $A$ به تکه‌های جدا از هم $A\cap B_1,\ A\cap B_2,\dots$ شکسته می‌شود. طبق قانون جمع، احتمال $A$ جمع احتمال این تکه‌هاست و طبق قانون ضرب $\P(A\cap B_i)=\P(B_i)\P(A\mid B_i)$. در نمودار درختی، این یعنی «جمع احتمال همهٔ مسیرهایی که به $A$ می‌رسند».

\begin{example}[label=hs:ex:factory]{کارخانه}
در یک کارخانه، دستگاه $M_1$ شصت درصد و دستگاه $M_2$ چهل درصد لامپ‌ها را تولید می‌کند. ۲٪ لامپ‌های $M_1$ و ۵٪ لامپ‌های $M_2$ معیوب‌اند. لامپی به تصادف انتخاب می‌شود. احتمال معیوب بودن آن چقدر است؟
\begin{solution}
$\P(\fa{معیوب})=0.6\times0.02+0.4\times0.05=0.012+0.020=0.032$؛ یعنی ۳٫۲٪.
\end{solution}
\end{example}

\begin{example}[label=hs:ex:bus]{دیر رسیدن سرویس}
احتمال بارانی بودن یک روز $0.3$ است. در روزهای بارانی سرویس مدرسه با احتمال $0.5$ و در روزهای غیربارانی با احتمال $0.1$ دیر می‌رسد. احتمال دیر رسیدن سرویس در یک روز دلخواه چقدر است؟
\begin{solution}
$\P(\fa{دیر})=0.3\times0.5+0.7\times0.1=0.15+0.07=0.22$.
\end{solution}
\end{example}

%% ------------------------------------------------------------
\section{قضیهٔ بیز: برگشت از نتیجه به علت}
%% ------------------------------------------------------------

در مثال دو جعبه پرسیدیم: «احتمال قرمز بودن مهره چقدر است؟» حالا سؤال را برعکس می‌کنیم: \key{مهره قرمز درآمده؛ احتمال اینکه از جعبهٔ ۱ برداشته شده باشد چقدر است؟} این یعنی برگشت از \emph{نتیجه} (قرمز) به \emph{علت} (کدام جعبه).

\begin{theorem}[label=hs:thm:bayes]{قضیهٔ بیز}
اگر $B_1,\dots,B_k$ افرازی از $S$ باشند و $\P(A)>0$، آنگاه برای هر $i$:
\[ \P(B_i\mid A)=\frac{\P(B_i)\,\P(A\mid B_i)}{\P(A)}, \]
که در آن مخرج، $\P(A)$، با قانون احتمال کل حساب می‌شود.
\end{theorem}

\textbf{چرا؟} طبق قانون ضرب، $\P(A\cap B_i)$ را به دو روش می‌نویسیم: $\P(B_i)\P(A\mid B_i)=\P(A)\P(B_i\mid A)$. با تقسیم بر $\P(A)$ قضیه به دست می‌آید.

\begin{intuition}
در نمودار درختی: احتمال اینکه از مسیر $B_i$ آمده باشیم، به شرط اینکه به $A$ رسیده‌ایم، برابر است با
\[ \frac{\text{\rl{احتمال مسیری که از $B_i$ به $A$ می‌رسد}}}{\text{\rl{جمع احتمال همهٔ مسیرهایی که به $A$ می‌رسند}}}. \]
\end{intuition}

\begin{example}[label=hs:ex:boxes-bayes]{دو جعبه، برعکس}
در مثال \ref{hs:ex:two-boxes} مهرهٔ برداشته‌شده قرمز است. احتمال اینکه از جعبهٔ ۱ باشد چقدر است؟
\begin{solution}
از نمودار درختی: مسیر «جعبهٔ ۱، قرمز» احتمال $\tfrac3{10}$ و کل مسیرهای قرمز $\tfrac25$ است. پس
\[ \P(\fa{جعبهٔ ۱}\mid\fa{قرمز})=\frac{3/10}{2/5}=\frac34. \]
قبل از دیدن مهره، احتمال جعبهٔ ۱ برابر $\tfrac12$ بود؛ قرمز بودن مهره (که در جعبهٔ ۱ رایج‌تر است) آن را به $\tfrac34$ رساند.
\end{solution}
\end{example}

\begin{example}[label=hs:ex:factory-bayes]{کدام دستگاه؟}
در مثال \ref{hs:ex:factory}، لامپ انتخاب‌شده معیوب است. احتمال اینکه دستگاه $M_2$ آن را ساخته باشد چقدر است؟
\begin{solution}
$\P(M_2\mid\fa{معیوب})=\dfrac{0.4\times0.05}{0.032}=\dfrac{0.020}{0.032}=0.625$. با اینکه $M_2$ فقط ۴۰٪ لامپ‌ها را می‌سازد، ۶۲٫۵٪ لامپ‌های معیوب مال آن است؛ چون درصد خرابی‌اش بیشتر است.
\end{solution}
\end{example}

\begin{example}[label=hs:ex:bus-bayes]{باران و سرویس}
در مثال \ref{hs:ex:bus}، سرویس امروز دیر آمده است. احتمال اینکه امروز بارانی باشد چقدر است؟
\begin{solution}
$\P(\fa{بارانی}\mid\fa{دیر})=\dfrac{0.3\times0.5}{0.22}=\dfrac{0.15}{0.22}\approx0.68$. دیر آمدن سرویس، احتمال باران را از ۳۰٪ به حدود ۶۸٪ رساند.
\end{solution}
\end{example}

\subsection{آزمایش پزشکی: یک نتیجهٔ شگفت‌انگیز}

\begin{example}[label=hs:ex:medical]{آزمایش ۹۵٪ دقیق}
یک بیماری در ۱٪ مردم وجود دارد. آزمایشی برای آن داریم که:
اگر فرد بیمار باشد، با احتمال ۹۵٪ جواب مثبت می‌دهد؛ و اگر فرد سالم باشد، با احتمال ۵٪ به اشتباه جواب مثبت می‌دهد.
آزمایش کسی مثبت شده است. احتمال اینکه واقعاً بیمار باشد چقدر است؟
\begin{solution}
بیشتر مردم حدس می‌زنند حدود ۹۵٪. بیایید با \emph{تعداد آدم‌ها} حساب کنیم. فرض کنید ۱۰۰۰۰ نفر آزمایش شوند:
\begin{figblock}
\begin{tikzpicture}[
  box/.style={draw=#1, fill=#1!10, rounded corners=3pt, minimum width=2.3cm, minimum height=0.8cm, font=\small, align=center},
  arr/.style={-{Stealth[length=2mm]}, thick, graytext}]
  \node[box=primary] (all) at (0,0) {\rl{۱۰۰۰۰ نفر}};
  \node[box=warncolor] (d) at (3.4,-1.5) {\rl{۱۰۰ بیمار}};
  \node[box=excolor] (h) at (-3.4,-1.5) {\rl{۹۹۰۰ سالم}};
  \node[box=warncolor] (dp) at (5,-3) {\rl{۹۵ مثبت}};
  \node[box=graytext] (dn) at (1.8,-3) {\rl{۵ منفی}};
  \node[box=accent] (hp) at (-1.8,-3) {\rl{۴۹۵ مثبت}};
  \node[box=graytext] (hn) at (-5,-3) {\rl{۹۴۰۵ منفی}};
  \draw[arr] (all) -- node[above right, font=\scriptsize]{\lr{$1\%$}} (d);
  \draw[arr] (all) -- node[above left, font=\scriptsize]{\lr{$99\%$}} (h);
  \draw[arr] (d) -- node[right, font=\scriptsize]{\lr{$95\%$}} (dp);
  \draw[arr] (d) -- (dn);
  \draw[arr] (h) -- node[right, font=\scriptsize]{\lr{$5\%$}} (hp);
  \draw[arr] (h) -- (hn);
\end{tikzpicture}

\figcaption[hs:fig:medical]{از ۱۰۰۰۰ نفر: ۹۵ بیمارِ مثبت و ۴۹۵ سالمِ مثبت.}
\end{figblock}
روی هم $95+495=590$ نفر جواب مثبت گرفته‌اند که فقط ۹۵ نفرشان واقعاً بیمارند:
\[ \P(\fa{بیمار}\mid\fa{مثبت})=\frac{95}{590}\approx0.16. \]
\textbf{فقط حدود ۱۶٪!} دلیلش این است که بیماری \emph{نادر} است: آن‌قدر آدم سالم زیاد است که حتی ۵٪ اشتباه در میان آن‌ها، از تعداد کل بیماران بیشتر می‌شود. با فرمول بیز هم همین به دست می‌آید:
\[ \P(\fa{بیمار}\mid\fa{مثبت})=\frac{0.01\times0.95}{0.01\times0.95+0.99\times0.05}=\frac{0.0095}{0.059}\approx0.16. \]
\end{solution}
\end{example}

\begin{warning}
$\P(\fa{مثبت}\mid\fa{بیمار})=0.95$ با $\P(\fa{بیمار}\mid\fa{مثبت})\approx0.16$ خیلی فرق دارد. اشتباه گرفتن این دو یکی از رایج‌ترین خطاها، حتی بین پزشکان، است. به همین دلیل پزشکان معمولاً آزمایش مثبت را با آزمایش دوم تأیید می‌کنند.
\end{warning}

\begin{moreinfo}{مسئلهٔ مونتی هال}
در یک مسابقه پشت یکی از سه در یک خودرو و پشت دو در دیگر بز است. شما در ۱ را انتخاب می‌کنید. مجری که می‌داند خودرو کجاست، یکی از دو در دیگر را که پشتش بز است باز می‌کند (مثلاً در ۳) و می‌پرسد: «می‌خواهید انتخابتان را به در ۲ عوض کنید؟»

جواب شگفت‌انگیز: \textbf{بله! با عوض کردن، احتمال برد از $\tfrac13$ به $\tfrac23$ می‌رسد.} دلیل ساده: انتخاب اول شما با احتمال $\tfrac23$ اشتباه است (خودرو پشت یکی از دو در دیگر است). در این حالت، مجری مجبور است تنها درِ بزدارِ باقی‌مانده را باز کند؛ پس در دیگری که باز نشده، حتماً خودرو دارد. یعنی با احتمال $\tfrac23$، عوض کردن شما را برنده می‌کند. (این را با قضیهٔ بیز هم می‌توان دقیق ثابت کرد.)
\end{moreinfo}

%% ------------------------------------------------------------
\section*{خلاصهٔ فصل ۶}
\addcontentsline{toc}{section}{خلاصه و تمرین‌ها}
%% ------------------------------------------------------------

\begin{summary}[title={\leavevmode\boxicon{accent}{$\Sigma$}\ \ خلاصهٔ فصل}]
\begin{itemize}
  \item \textbf{افراز:} پیشامدهای ناسازگاری که دقیقاً یکی‌شان رخ می‌دهد.
  \item \textbf{قانون احتمال کل:} $\P(A)=\sum_i\P(B_i)\P(A\mid B_i)$؛ در نمودار درختی: جمع مسیرهایی که به $A$ می‌رسند.
  \item \textbf{قضیهٔ بیز:} $\P(B_i\mid A)=\dfrac{\P(B_i)\P(A\mid B_i)}{\P(A)}$؛ برگشت از نتیجه به علت.
  \item برای مسائل بیز، فکر کردن با «تعداد آدم‌ها» (مثلاً ۱۰۰۰۰ نفر) خیلی کمک می‌کند.
  \item $\P(A\mid B)\ne\P(B\mid A)$؛ نادر بودن یک حالت (مثل یک بیماری) نتیجه را خیلی تغییر می‌دهد.
\end{itemize}
\end{summary}

\subsection*{تمرین‌ها}

\begin{exercise}
جعبهٔ ۱ دو مهرهٔ سفید و سه مهرهٔ سیاه و جعبهٔ ۲ چهار مهرهٔ سفید و یک مهرهٔ سیاه دارد. تاسی می‌اندازیم: اگر ۱ یا ۲ آمد از جعبهٔ ۱ و در غیر این صورت از جعبهٔ ۲ یک مهره برمی‌داریم. (الف) احتمال سفید بودن مهره چقدر است؟ (ب) اگر مهره سفید باشد، احتمال اینکه از جعبهٔ ۱ باشد چقدر است؟
\answer{(الف) $\tfrac13\cdot\tfrac25+\tfrac23\cdot\tfrac45=\tfrac{10}{15}=\tfrac23$ \quad (ب) $\tfrac{2/15}{2/3}=\tfrac15$}
\end{exercise}

\begin{exercise}
۶۰٪ دانش‌آموزان یک مدرسه دختر و ۴۰٪ پسرند. ۳۰٪ دخترها و ۵۰٪ پسرها عضو تیم ورزشی‌اند. (الف) چند درصد دانش‌آموزان عضو تیم‌اند؟ (ب) یک عضو تیم به تصادف انتخاب می‌شود؛ احتمال اینکه دختر باشد چقدر است؟
\answer{(الف) ۳۸٪ \quad (ب) $\tfrac{0.18}{0.38}=\tfrac9{19}\approx0.47$}
\end{exercise}

\begin{exercise}
سه دستگاه $A$، $B$ و $C$ به ترتیب ۵۰٪، ۳۰٪ و ۲۰٪ قطعات یک کارخانه را می‌سازند و درصد خرابی آن‌ها به ترتیب ۱٪، ۲٪ و ۳٪ است. (الف) احتمال معیوب بودن یک قطعه چقدر است؟ (ب) اگر قطعه‌ای معیوب باشد، احتمال اینکه دستگاه $C$ آن را ساخته باشد چقدر است؟
\answer{(الف) $0.017$ \quad (ب) $\tfrac{0.006}{0.017}=\tfrac6{17}\approx0.35$}
\end{exercise}

\begin{exercise}
یک بیماری در ۲٪ مردم وجود دارد. آزمایشی برای بیماران با احتمال ۹۰٪ و برای افراد سالم با احتمال ۱۰٪ جواب مثبت می‌دهد. با روش «۱۰۰۰ نفر» احتمال بیمار بودن کسی را که آزمایشش مثبت شده حساب کنید.
\answer{۱۸ بیمارِ مثبت و ۹۸ سالمِ مثبت: $\tfrac{18}{116}\approx0.16$}
\end{exercise}

\begin{exercise}
یک فیلتر ایمیل ۹۰٪ هرزنامه‌ها را تشخیص می‌دهد و ۲٪ ایمیل‌های معمولی را هم به اشتباه هرزنامه تشخیص می‌دهد. اگر ۲۰٪ ایمیل‌ها هرزنامه باشند، احتمال اینکه ایمیلی که فیلتر آن را هرزنامه تشخیص داده واقعاً هرزنامه باشد چقدر است؟
\answer{$\tfrac{0.2\times0.9}{0.2\times0.9+0.8\times0.02}=\tfrac{0.18}{0.196}\approx0.92$}
\end{exercise}


\end{document}
